\section{Shared results II: flag measures, row weights and metrics}\label{sec:pf:sh-flag}

\providecommand{\shflagmass}{\operatorname{mass}}
\providecommand{\shflagHS}[1]{\|#1\|_{\mathrm{HS}}}
\providecommand{\shflagcovN}{\operatorname{cov}_N}
\providecommand{\shflagvolR}{\operatorname{vol}_R}
\providecommand{\shflagII}{\mathrm{II}}
\providecommand{\shflagrel}{\mathcal T^{\mathrm{rel}}}

This section proves sixteen results of the formal development \cite{LeanSource2026} which are used by
several of the preceding sections of Part~II and which were cited there by their formal names. They fall
into six groups.
\begin{enumerate}[label=(\arabic*)]
\item \emph{Hermitian forms over CM fields} (\cref{sec:pf:sh-flag:herm}): the representation of totally
positive elements by binary hermitian forms (\cref{pf:sh-flag:hermrep}, used in \cref{sec:pf:multnc},
\cref{sec:pf:thmS} and \cref{sec:pf:top-shape}) and the ratio form of Siegel's mass formula for
orthonormal lattices (\cref{pf:sh-flag:massratio}, used in \cref{sec:pf:bundles} and
\cref{sec:pf:open0}).
\item \emph{Consequences of the Siegel--Weil formula} (\cref{sec:pf:sh-flag:sw}): the volume identities
$\sfname{SWXS}$ (\cref{pf:sh-flag:swxs}), the step-choice constants (\cref{pf:sh-flag:stepchoice}), the
chain of step estimates for equal words (\cref{pf:sh-flag:chain}), the non-generic row bound
$\sfname{RowNGP}$ (\cref{pf:sh-flag:rowng}), and Wu decay under the gap hypothesis
(\cref{pf:sh-flag:wu}); they are used in \cref{sec:pf:open0}, \cref{sec:pf:rows}, \cref{sec:pf:stageU}
and \cref{sec:pf:thmS}.
\item \emph{The compression formula for the lowest $K$-type} of a discrete series representation of
$\Uni(p,q)$ with large gap (\cref{pf:sh-flag:formulagap}, used in \cref{sec:pf:rows}).
\item \emph{Covolumes and amplitudes} (\cref{sec:pf:sh-flag:cov}): the covolume of the integral points of
the Bessel unipotent group (\cref{pf:sh-flag:covn}), the bound for the covolumes of the stabilisers of
relevant terms (\cref{pf:sh-flag:stabcov}) and the factorisation of the amplitudes of relevant terms into
local Bessel amplitudes (\cref{pf:sh-flag:jfact}); they are used in \cref{sec:pf:thmS}.
\item \emph{The type bound} (\cref{sec:pf:sh-flag:type}) for the degenerate relevant terms of the diagonal
orbital expansions (\cref{pf:sh-flag:typebound}, used in \cref{sec:pf:open0} and \cref{sec:pf:stageU}).
\item \emph{Intersection systems and small sections} (\cref{sec:pf:sh-flag:geo}): the construction of a
full arithmetic intersection system (\cref{pf:sh-flag:ais}), the packing bound for small sections
(\cref{pf:sh-flag:packing}), the integrability of point integrals (\cref{pf:sh-flag:ptsint}) and the
comparison of Chern-form integrals with generic intersection numbers (\cref{pf:sh-flag:deggen}); they are
used in \cref{sec:pf:deepstr}, \cref{sec:pf:tower} and \cref{sec:pf:towergeo}.
\end{enumerate}
Several of these results are reformulations or specialisations of other statements of the development, and
their proofs are short. Others (\cref{pf:sh-flag:typebound,pf:sh-flag:covn,pf:sh-flag:stabcov,%
pf:sh-flag:jfact,pf:sh-flag:formulagap,pf:sh-flag:ais,pf:sh-flag:packing}) are the top of long formal
arguments. For these we give the complete argument at the top levels and state the ingredients precisely;
the deeper ingredients are cited by their formal names (with a statement) and are listed in the report on
this section.

\paragraph{Conventions.}
We use the global data of \cref{sec:families}: a datum $D$ with ranks $m\ge2$, $k$, $f=[F:\Q]$,
$r_S=\lfloor m/2\rfloor-1$ and $m_R=2r_S+2$, its CM host $E/F$ with complex conjugation $x\mapsto\bar x$, the root discriminant
$\mathrm{rd}_E=|d_E|^{1/[E:\Q]}$ and $v=\mathrm{rd}_E/2$, a CM type $\Psi$ of $E$ (whose elements we denote
$\tau$; $|\Psi|=f$), size data $L_2$ with scale $b=e^L$, Hida settings $X$, full systems $Y$ over $X$, and
words $u,u'$. For words $u,u'$ we write $\mathcal T(u,u')$ for the term system of the orbital expansion
relating the column spaces of $u'$ and $u$, with the notation of \cref{pf:thmS:terms} and
\cref{pf:open0:def:standing}: terms $t$ with operators $\mathrm{op}(t)$ and vectors $\mathrm{vec}(t)$, the
weight $G(t)$, the sets $\mathrm{Deg}$, $\mathrm{NG}$, $\mathrm{Gen}$, the pair sums $\mathrm{ps}$, and
$d_u$, $D_*$, $dd(u,u')$; the pair $(X,Y)$ may be \emph{concretely tied} (\cref{pf:open0:def:standing}).
The \emph{$K_2$-hypotheses} $\mathrm{K}_2(C',C_l,C_g,c_b,C_w,c_T,C_T,C_0)$ for $(D,L_2)$ are the hypotheses
of the development with this name; they are implied by the standing hypotheses $\mathrm{H}_0(C_0)$ of
\cref{pf:open0:def:standing} (\lref{Formalism.MainNotions.NonGenericAndKTwoTheorems.Hyp_of}) and
contain, among others, the size inequalities $C'm^5fL\le k$ and the conditions on the host which we use
below. The \emph{counting context} $\mathrm{Ctx}(C_g,\dots,C_0,C',C_l;D,L_2,X,Y,u,u')$ is the conjunction
of the ties of $D$, the standing hypotheses $\mathrm{H}_0(C_0)$, the standard conditions for the pair
$(u,u')$, the concrete tie of $(X,Y)$ and the absolute convergence of $\mathcal T(u,u')$. Statements
``for all constants there are thresholds'' have the same shape as the uniform statements of
\cref{sec:pf:open0}: for all reals $C_g,c_b,C_w,c_T,C_T,C_0$ (with $c_b,c_T>0$ where indicated) there are
$C'_0,C_{l,0}$ such that the assertion holds for all $C'\ge C'_0$ and $C_l\ge C_{l,0}$ and all data in the
indicated hypotheses.

\begin{definition}[The Siegel--Weil formula]\label{pf:sh-flag:defsw}
The Siegel--Weil formula is used in the form of the development (\cref{lit:families:sw}): $(\mathrm{SW})$ is the
statement that for every CM extension $E/F$, every $M$ and every non-degenerate hermitian $J\in M_M(E)$ the
following two assertions hold for $H=\Uni(J)$.
\begin{itemize}
\item[$(\mathrm{SW}_1)$] There is a Haar measure on $H(\A_F)$ and a measurable fundamental domain for
$H(F)$, and for every Haar measure $\mu$ and every such fundamental domain $\mathcal F$,
$\mu(\mathcal F)<\infty$.
\item[$(\mathrm{SW}_2)$] For every $n$ with $2n<M$, every hermitian $T\in M_n(E)$ with $\det T\ne0$ and every
factorisable $\Phi=\Phi_\infty\otimes\Phi_{\mathrm f}$ on $(\A_E^M)^n$ with $\Phi_\infty$ smooth,
compactly supported and rapidly decreasing and $\Phi_{\mathrm f}=\bigotimes_v\Phi_v$ pure, for all choices of
Haar measures, fundamental domains, characters of $\A_F/F$ and self-dual measures, there are continuous
densities $g_\infty$, $g_v$ of the push-forwards of $\Phi_\infty$, $\Phi_v$ under the Gram map, the product
$\prod_vg_v(T)$ is multipliable, and the mean of $h\mapsto\sum_{X\in(E^M)^n,\ (h_J(x_a,x_b))=T}\Phi(h^{-1}X)$
over $\mathcal F$ equals $g_\infty(T)\prod_vg_v(T)$.
\end{itemize}
\formal{HeckeCharacterZeroLocus.SiegelWeilU.SiegelWeilUnitaryConvergent}
\end{definition}

%==============================================================================
\subsection{Hermitian forms over CM fields}\label{sec:pf:sh-flag:herm}

\begin{definition}\label{pf:sh-flag:defcmu}
Let $K$ be a number field with an involution $x\mapsto x^*$ (a field automorphism of order dividing $2$).
We say that $K$ satisfies $(\mathrm{CMU})$ if
\begin{enumerate}[label=(\roman*)]
\item $K$ is totally complex;
\item $\varphi(x^*)=\overline{\varphi(x)}$ for every embedding $\varphi\colon K\to\C$ and every $x\in K$;
\item for every maximal ideal $\mathfrak P$ of $\OO_K$ there is $\theta\in\OO_K$ with
$\theta-\theta^*\notin\mathfrak P$.
\end{enumerate}
\formal{Formalism.HermitianLatticesTheory.IsCMUnram}
\end{definition}

By (i) and (ii), $K$ is a CM field and $x\mapsto x^*$ is its complex conjugation
(\lref{GlobalTheory.OrthoLattice.isCMField_of_star}): the fixed field $K^+$ of the involution is totally
real, since $\varphi(x)=\varphi(x^*)=\overline{\varphi(x)}$ for $x\in K^+$; as $K$ is totally complex, some
$\varphi(x)$ is not real, so the involution is non-trivial and $[K:K^+]=2$; and complex conjugation of a CM
field is the unique automorphism inducing complex conjugation under every embedding. Condition (iii) says that
$K/K^+$ is unramified at all finite primes (\cref{pf:bundles:rem:mass}). An element $x\in K$ is
\emph{totally positive} if $\Re\varphi(x)>0$ for every embedding $\varphi\colon K\to\C$; if $x=x^*$, then
$\varphi(x)$ is real, so this means $\varphi(x)>0$.

For a CM field $E$ with complex conjugation $x\mapsto\bar x$ and $G\in M_r(E)$ put
$h_G(x,y)=\sum_{a,b}x_aG_{ab}\bar y_b$ for $x,y\in E^r$; $G$ is \emph{hermitian} if $\bar G_{ab}=G_{ba}$ for all
$a,b$.

\begin{littheorem}[Landherr {\cite{Landherr1936}}; see also {\cite[Ch.~10, \S6]{Scharlau1985}}]\label{pf:sh-flag:landherr}
Let $E$ be a CM field, $r\in\N$ and $G,G'\in M_r(E)$ hermitian with $\det G\ne0\ne\det G'$. Assume that for every
embedding $\tau\colon E\to\C$ there is $p_0$ such that both $\tau(G)$ and $\tau(G')$ are congruent over $\C$
(by matrices $Q$ with $\det Q\ne0$, in the sense $(h_{\tau(G)}(Qe_a,Qe_b))_{a,b}$) to
$\mathrm{diag}(1,\dots,1,-1,\dots,-1)$ with $p_0$ entries $1$, and that $\det G'=z\bar z\det G$ for some $z\in E$.
Then there is $\gamma\in\GL_r(E)$ with $h_G(\gamma x,\gamma y)=h_{G'}(x,y)$ for all $x,y\in E^r$.
\formal{GlobalTheory.Published.LandherrSt}
\end{littheorem}

This is Landherr's Hasse principle for hermitian forms: over a CM extension $E/E^+$, non-degenerate
hermitian forms of the same rank are isometric if they have the same signatures at the real places of $E^+$
and determinants in the same class of $(E^+)^\times/\Nm(E^\times)$. The development derives it from the Hasse
norm theorem for quadratic extensions (\lref{GlobalTheory.QNorm.landherr_and_landherrSt}).

\begin{proposition}[Totally positive elements are represented by binary forms]\label{pf:sh-flag:hermrep}
Let $K$ be a number field with an involution satisfying $(\mathrm{CMU})$, and let $a,b,n\in K$ be totally
positive with $a^*=a$, $b^*=b$ and $n^*=n$. Then there are $s,t\in K$ with
\[
a\,ss^*+b\,tt^*=n .
\]
\formallabel{GlobalTheory.hermRepUnram}\formallabel{GlobalTheory.T.binaryHermitianForm_represents_totallyPositive}
\formal{GlobalTheory.BinaryRep.represents_totallyPositive}
\end{proposition}

\begin{proof}
As explained above, $K$ is a CM field and $x^*=\bar x$ is its complex conjugation. Elements fixed by the
involution have real images under every embedding; since $a,b,n$ are totally positive, they are non-zero, and
$abn$ is fixed by the involution and totally positive. Consider the hermitian matrices
\[
G=\mathrm{diag}(a,b),\qquad G'=\mathrm{diag}(n,\,abn)\qquad\text{in }M_2(K).
\]
They are hermitian since their diagonal entries are fixed by conjugation, and $\det G=ab\ne0$,
$\det G'=n\cdot abn\ne0$. For an embedding $\tau$, both $\tau(G)$ and $\tau(G')$ are diagonal with positive real
entries $d_1,d_2$; with $Q=\mathrm{diag}(d_1^{-1/2},d_2^{-1/2})$ one gets
$(h_{\tau(G)}(Qe_a,Qe_b))_{a,b}=1$, so both are congruent to the identity ($p_0=2$). Finally
$\det G'=n\bar n\det G$ since $\bar n=n$. By \cref{pf:sh-flag:landherr} there is $\gamma\in\GL_2(K)$ with
$h_G(\gamma x,\gamma y)=h_{G'}(x,y)$. For $x=y=e_1$ and $(s,t)=\gamma e_1$ this reads
$a s\bar s+b t\bar t=h_{G'}(e_1,e_1)=n$.
\end{proof}

\begin{remark}\label{pf:sh-flag:rem:hermrep}
(a) Condition (iii) of $(\mathrm{CMU})$ is not used. (b) The development contains a second proof of the same
statement (\lref{GlobalTheory.Mass.binary_hermitian_represents}), which is the one used in the proof of
\cref{pf:sh-flag:massratio} below: with $w\ne0$ and $w^*=-w$, $\delta=w^2\in K^+$ is totally negative, and the
quinary quadratic form $ax_0^2-a\delta x_1^2+bx_2^2-b\delta x_3^2-nx_4^2$ over $K^+$ is indefinite at every real
place, hence isotropic by the Hasse--Minkowski theorem (derived in the development from the Hasse norm theorem);
a non-trivial zero with $x_4\ne0$ gives $s=(x_0+x_1w)/x_4$, $t=(x_2+x_3w)/x_4$.
\end{remark}

For a $K$-vector space $V$ with a hermitian form $h$ (sesquilinear, conjugate-linear in the first argument,
with $h(y,x)=h(x,y)^*$) and an $\OO_K$-submodule $L\subseteq V$, the \emph{mass} is
\[
\shflagmass(h,L)=\sum_{M\in\mathcal S}\frac1{\#\Aut(M)}
\]
if there is a finite set $\mathcal S$ of lattices in the genus of $L$ containing exactly one lattice of every
isometry class in the genus, and $\shflagmass(h,L)=0$ otherwise; $\Aut(M)$ is the group of isometries $g$ of $h$
with $gM=M$. An \emph{orthonormal lattice of rank $n$} for $h$ is the $\OO_K$-span of a $K$-basis
$e_1,\dots,e_n$ of $V$ with $h(e_i,e_j)=\delta_{ij}$. Let $d$ be the number of infinite places of $K$,
$v_K=|d_K|^{1/(2d)}/2$, and for $M\ge1$
\begin{gather*}
R_K(M)=\frac{\bigl((M-1)!\,(v_K/\pi)^M\bigr)^d}{\sqrt{(2v_K)^d}}\prod_{\mathfrak P}e_M(\mathfrak P),\\
e_M(\mathfrak P)=\begin{cases}\bigl(1-(-1)^M\Nm\mathfrak P^{-M/2}\bigr)^{-1},&\mathfrak P^*=\mathfrak P,\\
\bigl(1-\Nm\mathfrak P^{-M}\bigr)^{-1/2},&\mathfrak P^*\ne\mathfrak P,\end{cases}
\end{gather*}
the unconditional product running over the maximal ideals of $\OO_K$ (the right-hand side of
\eqref{pf:bundles:eq:mass}).

\begin{proposition}[Siegel mass ratio for orthonormal lattices]\label{pf:sh-flag:massratio}
Let $K$ satisfy $(\mathrm{CMU})$ and $M\ge3$. Let $h$, $h'$ be hermitian forms on $K$-vector spaces $V$, $V'$, and
let $L\subseteq V$, $L'\subseteq V'$ be orthonormal lattices of ranks $M$ and $M-1$. Then
\[
\frac{\shflagmass(h,L)}{\shflagmass(h',L')}=R_K(M).
\]
\formallabel{GlobalTheory.Mass.siegelMassRatio_holds}
\formal{Formalism.HermitianLattices.Ext.L3.siegelMassRatio_of_LI}
\end{proposition}

\begin{proof}
Let $F=K^+\subseteq K$ be the fixed field of the involution
(\lref{Formalism.HermitianLattices.Ext.L3.exists_fixF}). We use the following statement of the
development, \lref{Formalism.HermitianLattices.Ext.L3.massConst_all}: \emph{assume the representation
property of \cref{pf:sh-flag:hermrep} and the local density statement $\sfname{LIStat}$; then there is
$c_1>0$ such that for every $m\ge1$, every finite-dimensional $K$-vector space $W$ of dimension $m$ with a
hermitian form $k$ which is totally positive definite ($k(x,x)$ totally positive for $x\ne0$) and every
$\OO_K$-lattice $\Lambda\subseteq W$ which is integral for $k$,}
\begin{equation}\label{pf:sh-flag:eq:massconst}
\shflagmass(k,\Lambda)=c_1\prod_{j=2}^mR_K(j)\cdot\lambda_F(k,\Lambda),
\end{equation}
\emph{where $\lambda_F(k,\Lambda)=\prod_{\mathfrak p}\lambda_{\mathfrak p}(k,\Lambda)$ is the unconditional
product over the maximal ideals $\mathfrak p$ of $\OO_F$ of the ratios of the standard local density to the local
density of $\Lambda$ at $\mathfrak p$, normalised by the $\mathfrak p$-adic valuation of the discriminant.} It is
proved by induction on $m$, from the case $m=1$ and a step from $m-1$ to $m$
(\lref{Formalism.HermitianLattices.Ext.L3.massConst_succ}); $\sfname{LIStat}$ is a local identity for the
lattices obtained by splitting off a vector of non-zero norm. The two hypotheses hold: the representation property is
\cref{pf:sh-flag:hermrep} (in the formal proof, its second proof
\lref{GlobalTheory.Mass.binary_hermitian_represents}), and $\sfname{LIStat}$ is
\lref{Formalism.HermitianLattices.Ext.L3.liStat}.

An orthonormal lattice $L$ of rank $M$ spans $V$, which is therefore of dimension $M$; $h$ is totally positive
definite, since $h(x,x)=\sum_ix_ix_i^*$ in the coordinates of an orthonormal basis; $L$ is a lattice and
$h(L,L)\subseteq\OO_K$. Moreover $\lambda_F(h,L)=1$ for orthonormal lattices
(\lref{Formalism.HermitianLattices.Ext.L3.lamG_eq_one_of_ortho}). The same holds for $(h',L')$ in rank
$M-1\ge2$. Hence \eqref{pf:sh-flag:eq:massconst} gives
\[
\shflagmass(h,L)=c_1\prod_{j=2}^{M}R_K(j),\qquad \shflagmass(h',L')=c_1\prod_{j=2}^{M-1}R_K(j).
\]
Since $c_1>0$ and $R_K(j)>0$ for $j\ge2$ (\lref{Formalism.HermitianLattices.Ext.L3.massRatio_pos}), the
second mass is positive and the quotient is $R_K(M)$.
\end{proof}

\begin{remark}\label{pf:sh-flag:rem:massratio}
The proposition is the ratio form of Siegel's mass formula for unimodular hermitian lattices
\cite{Siegel1935,Prasad1989,GanHankeYu2001} used in \cref{prop:families:mass} (compare
\eqref{pf:bundles:eq:mass}). The
masses themselves are only determined up to the constant $c_1$ of \eqref{pf:sh-flag:eq:massconst}, which the
formal proof does not compute; only ratios are asserted.
\end{remark}

%==============================================================================
\subsection{Consequences of the Siegel--Weil formula}\label{sec:pf:sh-flag:sw}

\subsubsection{The volume identities and Wu decay}

\begin{definition}\label{pf:sh-flag:defswxs}
For a datum $D$ and a Hida setting $X$ for $D$, the automorphic instance of $X$ consists of a CM extension
$E/F$, a non-degenerate hermitian matrix $J\in M_n(E)$ and further data. $\sfname{SWXS}$ is the statement that
$(\mathrm{SW}_1)$ and $(\mathrm{SW}_2)$ hold for $(E/F,J)$, for every $D$ and every Hida setting $X$ for $D$
(\cref{prop:families:volratio}(ii)).
\formal{PI.VolStabS.MassDataS.SWXS}
\end{definition}

\begin{proposition}[\cref{prop:families:volratio}(ii)]\label{pf:sh-flag:swxs}
$(\mathrm{SW})$ implies $\sfname{SWXS}$.
\formallabel{GlobalTheory.St.swxs_of_siegelWeil}
\end{proposition}

\begin{proof}
Let $X$ be a Hida setting for a datum $D$. The extension $E/F$ of the instance of $X$ is a CM extension and its
matrix $J$ is hermitian and non-degenerate; both are part of the structure of the instance (the CM step and the
hermitian step of the mass data of $X$). Hence $(\mathrm{SW})$ applies to $(E/F,J)$; the conjunction of
$(\mathrm{SW}_1)$ and $(\mathrm{SW}_2)$ for $(E/F,J)$ is, verbatim, the statement $\sfname{SWXS}$ at $X$
(\lref{GlobalTheory.St.siegelWeil_volumes}, which records that $(\mathrm{SW})$ gives these two assertions for
every CM extension and every non-degenerate hermitian matrix).
\end{proof}

\begin{definition}[Gap hypothesis; Wu decay]\label{pf:sh-flag:defwu}
The gap hypothesis $\mathrm{GH}(D)$ for a datum $D$ of rank $r$ is: $r\ge1801$, and for every $\tau\in\Psi$ the
Harish-Chandra parameter $(\lambda_{\tau,j})_j$ of $D$ at $\tau$ satisfies
$\lambda_{\tau,j+1}-\lambda_{\tau,j}\ge16r^{19}$ for all consecutive indices. For a Hida setting $X$, let
$G_\infty=\prod_v\Uni(\eps_v)$ be the group of its archimedean points with its Haar measure $\mu_\infty$,
$\phi_X$ its archimedean test function, and for a word $u$ let $W_u$ be the archimedean Bessel coefficient of
$u$ and $\nu_u$ the quotient measure on $G_\infty$ attached to $u$. Put
$\|g\|_\infty=\prod_v\|g_v\|_{\mathrm{op}}$, $\varrho_u(x)=\inf_t\|tx\|_\infty$ ($t$ in the archimedean points of
the stabiliser group of $u$), and $N_1=4n^2\sqrt{2\max_vq_v}$, where $(p_v,q_v)$ are the signatures. \emph{Wu
decay} for $X$ is the statement that for every word $u$:
\begin{itemize}
\item (reproduction) there are a compact set $\Omega\ni1$, a bounded measurable function $\alpha$ on $G_\infty$
vanishing outside $\Omega$, such that $z\mapsto\alpha(z)W_u(xz)$ is $\mu_\infty$-integrable and
$\int\alpha(z)W_u(xz)\,d\mu_\infty(z)=W_u(x)$ for every $x$;
\item (weighted integrability) $\int\|W_u(x)\|\,\varrho_u(x)^{N_1}\,d\nu_u(x)<\infty$.
\end{itemize}
\formal{Formalism.L2.SG.ConcreteSettingExists.GapHyp}\ \formal{Formalism.L2.SG.WuDecay.WuDecay}
\end{definition}

\begin{proposition}[Wu decay with large gap]\label{pf:sh-flag:wu}
For every datum $D$ satisfying $\mathrm{GH}(D)$ and every Hida setting $X$ for $D$, Wu decay holds for $X$ and
$\phi_X$ is $\mu_\infty$-integrable.
\formallabel{GlobalTheory.wuDecayOfGap}
\end{proposition}

\begin{proof}
The development proves that this statement follows from the uniqueness of Bessel models in large gap
(\lref{Formalism.L2.SG.Pf.wuDecayOfGap_of_uniqGap}: if for all $k\ge1$ and $r$, all signatures and all
Harish-Chandra parameters with gap at least $120(2r+k+1)^3$, any two smooth Bessel maps of a discrete series
representation of $\Uni(2r+k+1)$, with respect to a generic character and a representation of the Bessel
subgroup, are linearly dependent, then Wu decay and integrability hold as stated). That statement is proved from
the uniqueness hypothesis, five further statements of the development (the non-vanishing of Bessel coefficients
in large gap, the weighted integrability of local coefficients, the existence of local coefficients, a product
structure of the quotient by the maximal compact subgroup, and the unimodularity of the local unipotent groups)
and a local scalar reduction. The uniqueness of Bessel models in large gap is
\lref{CompanionTheory.Bmu.besselMapUniqGap} (via \lref{CompanionTheory.Lime.BM.wuDecayOfGap}).
\end{proof}

\subsubsection{The step statements}

The non-generic row bound is assembled from a number of uniform \emph{step statements}. Each of them has the form
``for all constants there are thresholds such that for every $(D,L_2,X,Y,u,u')$ in the counting context
$\mathrm{Ctx}$ (and further data) a certain property holds''. Besides the counting context, they involve
\emph{mass-data ties} $(\beta,X_0,\psi)$ of $X$ (an isomorphism $\beta$ of the adelic bridge of the instance of $X$,
a frame $X_0\in(E^n)^{2r_S+2}$ and a character $\psi$ of $\A_F/F$, subject to the compatibilities required in
the development), \emph{mass data} built from such a tie and from parameters (three reals $s_d,C,\Theta$, a
flag $F_0$, a finite set of row lattices, a function on terms, a pair system, indices at the infinite places and
a family of measures on $G_\infty$), the \emph{admissibility} of mass data (with two reals $I_{\mathrm{rat}}$,
$I_{\mathrm{re}}$), and \emph{pinned} measures $\nu$ for $(u,u')$. We use the following ones.

\begin{definition}[Four step statements]\label{pf:sh-flag:defsteps}
\begin{itemize}
\item[$(\mathrm{SC})$] \emph{Step-choice constants}: for every context and every mass-data tie there are
parameters and reals $I_{\mathrm{rat}},I_{\mathrm{re}}$ such that the mass data are admissible.
\item[$(\mathrm{PE})$] \emph{Pin step}: for every context there is a pinned measure for $(u,u')$.
\item[$(\mathrm{FN})$] \emph{Form of non-generic terms}: for every context, every pinned $\nu$ and every non-generic
term $t$ of $\mathcal T(u,u')$ with $\mathrm{vec}(t)\ne0$, the point of $t$ is the image of a rational point $w$
having a real representative, and for every such $w$ the operator of $t$ has the form $c\,A\circ O_\nu(w)\circ B$
with $c\in\C$, operators $A$, $B$ on the archimedean spaces of $u$, $u'$, and $O_\nu(w)$ the archimedean orbital
integral of the Bessel coefficients $W_u,W_{u'}$ at the real point of $w$.
\item[$(\mathrm{TB})$] \emph{Pinned Theorem B}: for all constants and every $A\in\N$ there is $C_B\ge0$ such that,
beyond thresholds, for every context, mass-data tie, admissible mass data, pinned $\nu$ and auxiliary
measures, the point-sum input of the abstract non-generic analysis holds for $\mathcal T(u,u')$ with the
constants $A(r_S+1)^2$ and $\exp(C_B(2r_S+2)^5f\log b)\,I_{\mathrm{re}}$.
\end{itemize}
\formal{Formalism.L2.SG.NGAnalytic.StepChoiceConstantsC}\ \formal{Formalism.L2.SG.NGAnalytic.StepPinExists}\
\formal{Formalism.L2.SG.NGVanish.FormNGC}\ \formal{Formalism.L2.SG.NGAnalytic.StepThmBPinnedC}
\end{definition}

\begin{proposition}[Step-choice constants]\label{pf:sh-flag:stepchoice}
Assume $(\mathrm{SW})$ and the convergence at infinity of the real objects: for every datum $D$ with
$\mathrm{GH}(D)$, every Hida setting $X$ and all words $u,u'$, the real objects of $X$ at $(u,u')$ lie in
the convergence class at infinity with parameter $1$ (the hypothesis $\mathrm{InInfinityX}(X,u,u',1)$ of the
development). Then $(\mathrm{SC})$ holds.
\formallabel{GlobalTheory.St.stepChoiceDataC_of_siegelWeil_inInfinity}
\formallabel{Formalism.L2.SG.Pf.choice_massData}
\end{proposition}

\begin{proof}
By \cref{pf:sh-flag:swxs}, $\sfname{SWXS}$ holds. Given the constants, the size statements of the development
(\lref{Formalism.L2.SG.NGCut.sizes_of_hypotheses}) provide thresholds $s_0,t_0$; take
$C'_0=\max(s_0,29)$ and $C_{l,0}=\max(t_0,16)$. Let $C'\ge C'_0$, $C_l\ge C_{l,0}$, let $(D,L_2,X,Y,u,u')$ satisfy
the counting context and let $(\beta,X_0,\psi)$ be a mass-data tie. The size statements give $2v>0$ and the
positivity of the level of $X$, and the gap hypothesis $\mathrm{GH}(D)$ holds since $C'\ge29$ and $C_l\ge16$
(\lref{Formalism.L2.SG.ThmSOwn.gap_hypothesis_of_hypotheses}, applied to the standing hypotheses contained in
the context). Choose the trivial parameters: the reals
$(0,0,1)$, the flag $F_u$ of $u$, the set of row lattices of $u$ in $Y$, the zero function on terms, the trivial
pair system, the index $0$ at every infinite place and the zero family of measures. The admissibility of the
resulting mass data follows (\lref{Formalism.L2.SG.NGAnalytic.MassDataLemmas.massDataValid_bk_of}) from the
following conditions, each of which holds: the level, mass,
column and index conditions at the instance (statements of the development valid for every instance); the volume
condition, which is $\sfname{SWXS}$; the volume condition for the $\kappa$-part; the typed condition
(\lref{Formalism.L2.SG.Pf.MassDataX.nTy_holds}); the convergence of the products over the places
(\lref{Formalism.L2.SG.Pf.ProdConvCut2.prodConv_onTied}); the primitivity of the local frames at $u$ and at
$u'$; the $\kappa$-condition; the membership of the real objects in the class $\nu$; their convergence at infinity,
which is the hypothesis (applicable by $\mathrm{GH}(D)$); the positivity of the mass data
(\lref{Formalism.L2.SG.Pf.MassDataX.nPos_of}, from $\sfname{SWXS}$ and the product convergence); and the frame
volume condition (\lref{Formalism.L2.SG.Pf.MassDataX.frameVolAt_of_prodConv}). This is
\lref{Formalism.L2.SG.Pf.stepChoiceDataC_of_leavesG}.
\end{proof}

\begin{proposition}[The chain of step estimates for equal words]\label{pf:sh-flag:chain}
The following nine statements of the development hold.
\begin{enumerate}[label=(\arabic*)]
\item The \emph{non-generic chain}: the conjunction of thirteen abstract statements on term systems (exact
cores, sums of class squares, rows from classes, a Cauchy--Schwarz inequality for classes, two bracket
inequalities, a joint mean-amplitude bound, the abstract non-generic row, a norm bound, the frame passage, a real
part estimate, an exponent sum and the numerical non-generic inequalities).
\item The \emph{class step}: for every datum all of whose type primes are $\ge3$, every concretely tied $(X,Y)$
and every word $u$, the relation $\mathrm{scal}$ is an equivalence relation on the terms of $\mathcal T(u,u)$.
\item The \emph{point step}: for every context with $u'=u$, the point input of the abstract analysis holds for the
generic part of $\mathcal T(u,u)$.
\item The \emph{norm-one step}: for every context with $u'=u$, the two norms $N_G(\gamma)$ and $N_D(\gamma)$ of every
point $\gamma$ of the generic part of $\mathcal T(u,u)$ (normalised by the root discriminant, and for $N_G$ by
$T_u^2$) are at least $1$.
\item The \emph{tie step}: every context admits a mass-data tie.
\item The pin step $(\mathrm{PE})$.
\item The \emph{vanishing step}: for every context with $u'=u$ and every pinned $\nu$, a generic term whose pinned
local amplitude at its point vanishes has $\mathrm{vec}(t)=0$.
\item The form of non-generic terms $(\mathrm{FN})$.
\item The generic family step implies the step of local generation for pinned data.
\end{enumerate}
\formallabel{GlobalTheory.T.equalWords_nonGeneric_estimates_chain}
\formallabel{GlobalTheory.closed_steps_equal_words}
\end{proposition}

\begin{proof}
(1) is \lref{Formalism.FourierDecayAlt.P.norm_chain}, and (2) is
\lref{Formalism.L2.SG.Pf.stepClassG_holds}. (3) and (4) follow from the same ten statements of the
development, each valid for all data (among them the unimodularity of $G_\infty$, the triviality of the level
character, the rational cocompactness of the stabilisers and the non-vanishing of the leading minors at some
point), by \lref{Formalism.L2.SG.Pf.MainEqPoints.stepPointsG_of_leaves} and
\lref{Formalism.L2.SG.NGCut.LeadingMinorIntegrality.Gen.stepNormOneG_of_leaves}. (5) follows from the Gram vectors
of the closed steps (\lref{GlobalTheory.closed_steps_gram_vectors}). (6) is
\lref{Formalism.L2.SG.NGPin.stepPinExists}. (7) follows from the form of generic terms
(\lref{Formalism.L2.SG.Pf.formGC_holds}) by \lref{Formalism.L2.SG.Pf.MainEqVanish.stepTermVanishG_of_form}.
(8) is \lref{Formalism.L2.SG.Pf.formNGC_holds}. For (9), the local non-generic bound of the development
(\lref{GlobalTheory.local_nongeneric_bound_record_family} with the parameter $39$) and the generic family
step give the local step for pinned data (\lref{Formalism.L2.SG.Pf.stepLocPinnedG_of_fam}), and this gives the
step of local generation (\lref{Formalism.L2.SG.Pf.stepLocGenPinnedG_of_loc}).
\end{proof}

\subsubsection{The non-generic row}

\begin{theorem}[\cref{prop:families:volratio}(iii)]\label{pf:sh-flag:rowng}
Assume $(\mathrm{SW})$. Then $\sfname{RowNGP}$ holds: for all $C_g,c_b,C_w,c_T,C_T,C_0$ with $c_b,c_T>0$ there are
$C_4\ge0$ and $C'_0,C_{l,0}$ such that for all $C'\ge C'_0$, $C_l\ge C_{l,0}$, every datum $D$ with size data
$L_2$, every Hida setting $X$, every full system $Y$ and all words $u,u'$: if the ties of $D$, the standing
hypotheses $\mathrm{H}_0(C_0)$ (with $C',C_l,C_g,c_b,C_w,c_T,C_T$), the standard conditions for $(u,u')$ and the
concrete tie of $(X,Y)$ hold and $\mathcal T(u,u')$ is absolutely convergent, then
\[
\mathrm{ps}\bigl(\mathrm{one}(\mathrm{NG})\bigr)\le\eps_{\mathrm{NG}}(C_4)\,dd(u,u'),\qquad
\eps_{\mathrm{NG}}(C_4)=2^{-(k-2r_S-5)}\exp\bigl(C_4m_R^5f\log b\bigr),
\]
for the term system $\mathcal T(u,u')$.
\formallabel{GlobalTheory.rowNGP_of_siegelWeil}
\formallabel{GlobalTheory.rowNGP_of_siegelWeil_of_stepChoice}
\formallabel{GlobalTheory.T.corankOneNonGenericRowBound_closed}
\formallabel{GlobalTheory.NGRowClose.close}
\formallabel{GlobalTheory.NGRowClose.rowNGP_of_formNGVanish}
\end{theorem}

\begin{proof}
\emph{Step 1: the inputs.} By \cref{pf:sh-flag:swxs}, $\sfname{SWXS}$ holds. The convergence at infinity of the
real objects under the gap hypothesis holds by \lref{Formalism.L2.SG.RealObj.EnergyConvJ.BM.realObj_inInfinityX_jsz}
(applied with the uniqueness of Bessel models in large gap, \lref{CompanionTheory.Lime.BM.uniqGap_holds}), so
$(\mathrm{SC})$ holds by \cref{pf:sh-flag:stepchoice}. $(\mathrm{PE})$ and $(\mathrm{FN})$ hold
(\lref{Formalism.L2.SG.NGPin.stepPinExists} and \lref{Formalism.L2.SG.Pf.formNGC_holds}; these are (6) and (8) of
\cref{pf:sh-flag:chain}). $(\mathrm{TB})$
holds by \lref{Formalism.MainStatements.RowRest.analytic_pinnedTheoremB_uniform}, which deduces it from four
uniform statements of the development: a rotated joint norm bound (the uniform form of Theorem~B), a uniform
amplitude bound, the uniformity of points, and uniform row data (from the Bessel--Whittaker statement and a
uniform restriction bound). The closed row theorem below uses $(\mathrm{TB})$ in a weaker form, with an
additional hypothesis which is not used; this weaker form follows trivially.

\emph{Step 2: the four family statements.} The non-generic terms of corank one of $\mathcal T(u,u')$ are cut into
local families by their defect (\lref{GlobalTheory.T.non_generic_corank_one_terms_family_cut_by_defect}, which
identifies the classes of non-generic terms with the rows in a double coset). From this cut,
\lref{GlobalTheory.T.nongeneric_row_statements_of_family_statements} deduces four uniform statements on these
families from four statements of the development: (i) an orbit-family identity for non-generic families, from the
unfolding identity for the row-weight sums of families
(\lref{GlobalTheory.T.familyRowWeightSumUnfoldingIdentity}); (ii) the weight bound
$0<w_{\mathrm{tot}}\le c(2v,r_S,k,f)\,A\,2\,H(\gamma)^{e_N}n(\gamma)$ for the non-generic total weight of a family over
a point $\gamma$, in terms of the amplitude constant $A$ of the mass data, from the positivity and boundedness of
the generic masses (\lref{GlobalTheory.T.familyGenericMassPositiveBounded_proof}, which uses $\sfname{SWXS}$);
(iii) a product structure: the coefficients of the point sums are products of family amplitudes, after a
rescaling $s$ of the pinned measure with $s=1$ on regular points and $s^2\le b^{C_{\mathrm{St}}(2r_S+2)^3f}$, from
the factorisation of group sums into row weights and amplitudes
(\lref{GlobalTheory.T.familyGroupSumFactorsRowWeightAmplitude_proof}); (iv) the existence of local objects at
the terms (\lref{GlobalTheory.T.familyLocalObjectsExistAtTerms_proof}).

\emph{Step 3: assembly.} The bound now follows from $\sfname{SWXS}$, $(\mathrm{SC})$, $(\mathrm{PE})$, $(\mathrm{FN})$,
the weak form of $(\mathrm{TB})$ and the four family statements by
\lref{GlobalTheory.T.corankOnePairsNonGenericRowBound_of_remaining_premises}. In that theorem, the closed steps
of the development (classes, normal invariants, points, norm one and Gram vectors,
\lref{GlobalTheory.RowNGClosed.closedSteps}) and the family statements give the family step
(\lref{GlobalTheory.T.familyStep_of_four_globalLeaves}); with the local non-generic bound at the parameter
$39$ it gives the step of local generation for pinned data; $(\mathrm{FN})$ gives the vanishing step for pinned
data; the pin, vanishing, local and Theorem-B steps give the analytic facts of the row
(\lref{GlobalTheory.T.rowAnalyticFacts_of_four_steps}); with $(\mathrm{SC})$ and a mass-data tie they give
admissible mass data with analytic constants (\lref{GlobalTheory.T.existsRowLevelWithAnalyticFacts_of_steps}),
and hence a splitting of the row into finite and archimedean parts
(\lref{GlobalTheory.T.rowSplitFiniteArchimedeanParts_of_steps}, using the norm-one step for the heights of the
points). Finally \lref{GlobalTheory.T.nonGenericCorankOnePairsBound_of_splitFiniteRealRow} applies the abstract
joint non-generic row bound of the development (\lref{Formalism.FourierDecayAlt.Final.pn_rowNGJoint}) and the
non-generic chain to the split row; the abstract bound has the form $\eps_{\mathrm{NG}}(C_4)\,dd(u,u')$ with a
constant $C_4\ge0$ depending only on the constants.
\end{proof}

\begin{remark}\label{pf:sh-flag:rem:rowng}
Of $(\mathrm{SW})$ the proof uses only $\sfname{SWXS}$, through $(\mathrm{SC})$ and the generic masses; the bound is
uniform in the words $u,u'$ (equal or not).
\end{remark}

%==============================================================================
\subsection{The compression formula for the lowest $K$-type}\label{sec:pf:sh-flag:lowk}

Let $n\ge1$ and $\eps\colon\{1,\dots,n\}\to\{\pm1\}$, $J_\eps=\mathrm{diag}(\eps_1,\dots,\eps_n)$, $G_\eps=\{g\in
\GL_n(\C):g^*J_\eps g=J_\eps\}\cong\Uni(p,q)$, $K_\eps=\{k\in\Uni(n):kJ_\eps=J_\eps k\}\cong\Uni(p)\times\Uni(q)$ with its
Haar probability measure $dk$, $\mathfrak g_\eps=\{X:X^*J_\eps+J_\eps X=0\}$ and
$\mathfrak p_\eps=\{X\in\mathfrak g_\eps:XJ_\eps=-J_\eps X\}$; let $r=\min(p,q)$ be the real rank. For $X\in\mathfrak
p_\eps$ the matrix $\exp X$ is hermitian positive definite, and so are its conjugates $k^{-1}\exp(X)k$, $k\in K_\eps$.

A \emph{Harish-Chandra parameter} is a strictly decreasing $\lambda\in\R^n$ with $\lambda_a-\frac{n-1}2\in\Z$ for all
$a$; it has \emph{gap at least $h$} if $|\lambda_a-\lambda_b|\ge h$ whenever $\eps_a\ne\eps_b$. Let $\rho_c$, $\rho_n$ be the
half-sums of the positive compact and non-compact roots (in the coordinates of the diagonal torus, for the ordering
of the indices), and $\Lambda_\lambda=\lambda+\rho_n-\rho_c$. A \emph{discrete series representation with parameter
$\lambda$} is a continuous unitary representation $\sigma$ of $G_\eps$ on a Hilbert space $H\ne0$ which is irreducible,
has square-integrable matrix coefficients, has a $K_\eps$-type of highest weight $\Lambda_\lambda$, and such that
$\|\Lambda_\lambda+2\rho_c\|\le\|\mu+2\rho_c\|$ for the highest weights $\mu$ of all its $K_\eps$-types. Let
$S\subseteq H$ be the sum of the $K_\eps$-types of $\sigma$ of highest weight $\Lambda_\lambda$ and $d_S=\dim S$. For a
matrix $z$ whose leading principal minors $\Delta_1(z),\dots,\Delta_n(z)$ are non-zero put
\[
\Phi_\lambda(z)=\prod_{i=1}^n\Delta_i(z)^{-s_i},\qquad s_i=\mu_i-\mu_{i+1},\quad \mu_i=\lambda_i+\tfrac{n+1}2-i\in\Z,\
\mu_{n+1}=0 .
\]

\begin{theorem}[Compression formula for the lowest $K$-type]\label{pf:sh-flag:formulagap}
Let $\lambda$ be a Harish-Chandra parameter with gap at least $h\ge120\,n\,r^{3/2}$, $\sigma$ a discrete series
representation of $G_\eps$ with parameter $\lambda$, and $w\in S$ a unit vector of weight $\Lambda_\lambda$ for the
diagonal torus of $K_\eps$. Then for every $X\in\mathfrak p_\eps$ and all $u,v\in S$,
\begin{equation}\label{pf:sh-flag:eq:compression}
\langle v,\sigma(\exp X)u\rangle=d_S\int_{K_\eps}\Phi_\lambda\bigl(k^{-1}\exp(X)\,k\bigr)\,
\langle v,\sigma(k)w\rangle\,\langle\sigma(k)w,u\rangle\,dk .
\end{equation}
\formallabel{GlobalTheory.Jet.FormGapFree.formulaGap}
\end{theorem}

Since every element of $G_\eps$ is of the form $k_1\exp(X)k_2$, the formula determines the matrix coefficients of the
lowest $K$-type; the decay estimates \ref{pf:rows:S1} of \cref{sec:pf:rows} are derived from it. The proof compares
both sides as functions on $G_\eps$: they satisfy the same first-order system (the equations of Schmid expressing that
the lowest $K$-type is annihilated by the lowering part of $\mathfrak p_\eps$), they agree to infinite order at $1$ by an
algebraic uniqueness statement for jets, and they are real-analytic along one-parameter subgroups. We use the
following notation. For a function $f$ on $G_\eps$ which is smooth near a point, and $1\le a,b\le n$, let $E_{ab}f$ be the
derivative of $f$ along the element of $\mathfrak g_\eps\otimes\C$ corresponding to the matrix unit $e_{ab}$, and for a word
$\omega=((a_1,b_1),\dots,(a_l,b_l))$ let $E_\omega f=E_{a_1b_1}\cdots E_{a_lb_l}f$.

\begin{lemma}[Analyticity of the lowest $K$-type]\label{pf:sh-flag:lowkan}
Let $\sigma$ be a continuous unitary representation of $G_\eps$ with a space $V$ of smooth vectors and its derivative
action $E_{ab}$ on $V$, and assume that the sum $S$ of the $K_\eps$-types of $\sigma$ of highest weight $\Lambda_\lambda$ is a single
$K_\eps$-type and is contained in $V$. Then for every
$x\in S$ and every $X\in\mathfrak g_\eps$, the map $t\mapsto\sigma(\exp(tX))x$ is real-analytic on $\R$.
\formallabel{GlobalTheory.Jet.Anal.lowK_analytic}
\end{lemma}

\begin{proof}
The $K$-type $S$ is the span, under the derivative action of $\mathfrak k_\eps\otimes\C$, of a highest weight vector $v_0$
(\lref{GlobalTheory.Jet.KType.kType_iff}), and $v_0$ is an eigenvector of the Casimir operator
(\lref{GlobalTheory.Jet.Anal.cas_eigen}). From this the development deduces a growth bound for the powers of
$d\sigma(X)$ on $S$ (\lref{GlobalTheory.Jet.Anal.pow_growth}): there are $B,C\ge0$ with
$\|d\sigma(X)^kx\|\le C^k(B+k)^k\|x\|$ for all $x\in S$ and $k\ge0$. Let $f(t)=\sigma(\exp(tX))x$. Then $f$ is infinitely
differentiable with $f^{(k)}(t)=\sigma(\exp(tX))\,d\sigma(X)^kx$, and by unitarity $\|f^{(k)}(t)\|\le M_k:=C^k(B+k)^k\|x\|$
for all $t$. For $\rho=1/(Ce+1)$ we have $\rho Ce<1$, and since $(B+k)^k\le e^Bk^k\le e^{B+k}k!$,
\[
\sum_k\frac{\rho^kM_k}{k!}\le e^B\|x\|\sum_k(\rho Ce)^k<\infty .
\]
By Taylor's formula with the remainder bounded by $\rho^{k+1}M_{k+1}/(k+1)!$, the Taylor series of $f$ at every $t_0$
converges to $f$ on $|t-t_0|<\rho$ (\lref{GlobalTheory.Jet.Anal.analytic_chain}); hence $f$ is analytic on $\R$.
\end{proof}

\begin{lemma}[Vanishing from zero jets]\label{pf:sh-flag:zerojet}
Let $f$ be a function on $G_\eps$ which is smooth near every point of $G_\eps$, and $X\in\mathfrak g_\eps$. If $E_\omega f(1)=0$
for every word $\omega$ and $t\mapsto f(\exp(tX))$ is real-analytic on $\R$, then $f(\exp X)=0$.
\formallabel{GlobalTheory.Jet.Line.vanish_of_zero_jet}
\end{lemma}

\begin{proof}
The $k$-th derivative of $t\mapsto f(\exp(tX))$ at $t=0$ is $(\tilde X^kf)(1)$ for the left-invariant vector field $\tilde X$;
writing $X$ in terms of the $E_{ab}$, it is a linear combination of the values $E_\omega f(1)$ for words of length $k$
(\lref{GlobalTheory.Jet.Line.iteratedDeriv_line}), hence $0$. The analytic function $t\mapsto f(\exp(tX))$ therefore
vanishes to infinite order at $0$, hence identically on the connected set $\R$; take $t=1$.
\end{proof}

The derivative action and jets are related to the following space of ``lowering'' vectors. Let $\Lambda=\Lambda_\lambda$,
let $V$ and $E_{ab}$ be as in \cref{pf:sh-flag:lowkan}, $u\in V$ a highest weight vector of weight $\Lambda$ for the compact
part, with finite-dimensional $K$-span $S_u$. Let $V^-(u)$ be the sum of the isotypic components, for the weights
$\Lambda+e_a-e_b$ with $b<a$ and $\eps_a\ne\eps_b$, of the space of arrays $(G_{ab})_{a,b}$ with $G_{ab}\in S_u$ (with the
action of $\mathfrak k_\eps$ on the array combining the action on the entries and on the indices). By the theorem of Schmid
on the lowest $K$-type \cite{Schmid1975}, in the form proved in the development
(\lref{GlobalTheory.Jet.Schmid.mul_vminus_eq_zero}), $\sum_{a,b}E_{ab}G_{ab}=0$ for every $G\in V^-(u)$ if $\sigma$ is a
discrete series representation with parameter $\lambda$ and $u$ is a highest weight vector of its lowest $K$-type.

\begin{lemma}[Zero jets]\label{pf:sh-flag:germ}
Let $\sigma$, $V$, $E_{ab}$, $\lambda$ and $u$ be as above, and assume that there is $g_0>2n$ with $\lambda_a-\lambda_b\ge g_0$
whenever $a<b$ and $\eps_a\ne\eps_b$. Let $F$ be a function from $G_\eps$ to the operators on $H$ such that for all $g\in
G_\eps$, $x\in S_u$, $k_1,k_2\in K_\eps$ and $y\in H$:
\begin{enumerate}[label=(\alph*)]
\item $F(g)x\in S_u$ and $F(k_1gk_2)x=\sigma(k_1)F(g)\sigma(k_2)x$;
\item $g\mapsto\langle y,F(g)x\rangle$ is smooth;
\item for every $G\in V^-(u)$, $\sum_{a,b}E_{ab}\bigl(g\mapsto\langle y,F(g)G_{ab}\rangle\bigr)=0$ near $1$;
\item $F(1)x=0$.
\end{enumerate}
Then $E_\omega\bigl(g\mapsto\langle y,F(g)x\rangle\bigr)(1)=0$ for every word $\omega$, every $x\in S_u$ and every $y\in H$.
\formallabel{GlobalTheory.Jet.Germ.germ_jet_zero}
\end{lemma}

\begin{proof}
By (a) and (b) there is a family of vectors $J(x,\omega)\in V$, linear in $x\in S_u$, representing the jets:
$\langle y,J(x,\omega)\rangle=E_\omega(\langle y,F(\cdot)x\rangle)(1)$ for all $y$
(\lref{GlobalTheory.Jet.Germ.exists_jet}). These jets satisfy: (i) $J(x,(c,d)\omega)=E_{cd}J(x,\omega)$ if $\eps_c=\eps_d$,
by the left equivariance in (a) (compact derivatives at $1$ act through $\sigma$); (ii) the commutation relations
$J(x,\omega(a,b)(c,d)\omega')-J(x,\omega(c,d)(a,b)\omega')=\delta_{bc}J(x,\omega(a,d)\omega')-\delta_{da}J(x,\omega(c,b)\omega')$ of
$\mathfrak{gl}_n$; (iii) $J(E_{cd}x,\omega)=J(x,\omega(c,d))$ if $\eps_c=\eps_d$, by the right equivariance in (a); (iv)
$\sum_{a,b}J(G_{ab},\omega(a,b))=0$ for $G\in V^-(u)$, by (c); and (v) $J(x,\varnothing)=0$, by (d).

These relations say that the jets form a module for $\mathfrak{gl}_n$ (\lref{GlobalTheory.Jet.Germ.jets_vanish}). Let
$\mathcal J$ be the space of families $(J(\omega))_\omega$ of vectors of $V$ with $J(\varnothing)\in S_u$ which satisfy the analogues of
(i) and (ii). On $\mathcal J$ the shifts $(E'_{cd}J)(\omega)=J(\omega(c,d))$ satisfy the relations of $\mathfrak{gl}_n$; by (iii),
$\iota\colon x\mapsto J(x,\cdot)$ is a map $S_u\to\mathcal J$ compatible with the compact derivatives; evaluation at the empty word,
$\mathrm{ev}(J)=J(\varnothing)$, takes values in $S_u$ and is compatible with the compact derivatives by (i). By (iv)
and (v), Schmid's relations $\sum_{a,b}E'_{ab}\iota(G_{ab})=0$ ($G\in V^-(u)$) and $\mathrm{ev}\circ\iota=0$ hold. The algebraic uniqueness
statement of the development (\lref{GlobalTheory.T.jet_uniqueness_algebraic}) now applies: \emph{for a $\mathfrak{gl}_n$-module $V'$
and maps $\iota\colon S_u\to V'$, $\mathrm{ev}\colon V'\to S_u$ compatible with the compact parts of the actions, Schmid's relations for
$\iota$ and $\mathrm{ev}\circ\iota=0$ imply $\mathrm{ev}(E'_\omega\iota(x))=0$ for every word $\omega$ and every $x\in S_u$, provided
$\lambda_a-\lambda_b\ge g_0>2n$ whenever $a<b$ and $\eps_a\ne\eps_b$.} Since $\mathrm{ev}(E'_\omega\iota(x))=J(x,\omega)$, all jets vanish.
\end{proof}

\begin{proof}[Proof of \cref{pf:sh-flag:formulagap}]
\emph{Step 1: the lowest $K$-type.} Since $\lambda$ is strictly decreasing with gap at least $h\ge120nr^{3/2}$, there is
$g_0>2n$ with $\lambda_a-\lambda_b\ge g_0$ whenever $a<b$ and $\eps_a\ne\eps_b$
(\lref{GlobalTheory.Jet.FormGap.exists_gap}). As the gaps exceed $2n$, $S$ is a single $K_\eps$-type of highest
weight $\Lambda_\lambda$ (\lref{GlobalTheory.Jet.MultGap.isKTypeOf_lowKSub}). There are a space $V\supseteq S$ of smooth
vectors with a derivative action (\lref{GlobalTheory.DerEx.exists_isDer_lowKSub}); $S$ is the $K$-span of $w$, which
is a highest weight vector, and $S$ is finite-dimensional (\lref{GlobalTheory.Jet.KType.kType_iff}). Let
$K_{\eps,\C}$ be the group of invertible matrices commuting with $J_\eps$. There is a representation $\tau$ of $K_{\eps,\C}$ on
$S$ with holomorphic matrix entries, irreducible under $K_\eps$, with $\tau(y)^*=\tau(y^*)$, $\tau(k)x=\sigma(\bar k)x$ for
$k\in K_\eps$ and $x\in S$, such that $w$ has weight $-\Lambda_\lambda$ for $\tau$ and is fixed by the lower triangular unipotent
elements of $K_{\eps,\C}$ (\lref{GlobalTheory.Dec.TauExt.tauExists}).

\emph{Step 2: the candidate.} Let $U_1$ be the set of matrices $z$ such that $k^{-1}zk$ has non-zero leading principal
minors for every $k\in K_\eps$; it is open and contains the positive definite matrices. For $z\in U_1$ put
\[
\psi(z)=d_S\int_{K_\eps}\Phi_\lambda(k^{-1}zk)\,\langle\tau(k)w,\cdot\rangle\,\tau(k)w\,dk\quad\in\End(S),
\]
and for an invertible $g$ let $|g|=(g^*g)^{1/2}$ and $\Psi(g)=\tau(g|g|^{-1})\,\psi(|g|)$. Then $\Psi(1)=\mathrm{id}_S$,
$\Psi(k\exp Y)=\tau(k)\psi(\exp Y)$ for $k\in K_\eps$ and $Y\in\mathfrak p_\eps$, and $\Psi(k_1gk_2)=\tau(k_1)\Psi(g)\tau(k_2)$ for
$g\in G_\eps$ and $k_1,k_2\in K_\eps$ (\lref{GlobalTheory.Dec.psiPolar}); the matrix entries of $\Psi$ are real-analytic
near every point of $G_\eps$ (\lref{GlobalTheory.Dec.psiAnalytic}). Writing $P_S$ for the orthogonal projection onto
$S$ and using $\tau(k)w=\sigma(\bar k)w$ and the invariance of $dk$ under $k\mapsto\bar k$, one has
\begin{equation}\label{pf:sh-flag:eq:psi}
\langle P_Sy,\psi(\bar z)P_Sx\rangle=d_S\int_{K_\eps}\Phi_\lambda\bigl(\overline{k^{-1}zk}\bigr)\,\langle y,\sigma(k)w\rangle\,
\langle\sigma(k)w,x\rangle\,dk\qquad(\bar z\in U_1)
\end{equation}
(\lref{GlobalTheory.Jet.FormGap.ps_entry}); near $1$ the same holds with $\Psi$ in place of $\psi$, also after
applying the derivatives $E_{ab}$ to both sides (\lref{GlobalTheory.Jet.FormGap.psi_near_one}). Consider the difference
\[
F(z)=P_S\,\sigma(z)-\Psi(\bar z)\,P_S
\]
(the second term viewed as an operator on $H$ with values in $S$).

\emph{Step 3: the hypotheses of \cref{pf:sh-flag:germ} for $F$ and $u=w$.} (a) $F(g)$ maps $S$ into $S$, and
$F(k_1gk_2)=\sigma(k_1)F(g)\sigma(k_2)$ on $S$, by the equivariance of $\sigma$, $P_S$ and $\Psi$ and the relation
$\tau(\bar k)=\sigma(k)$ on $S$ (\lref{GlobalTheory.Jet.FormGap.diffOp_equiv}). (b) The matrix entries of
$P_S\sigma(\cdot)x$ are smooth since $x$ is a smooth vector, and those of $\Psi$ are analytic. (c) Let $G\in V^-(w)$. Near
$1$, by \lref{GlobalTheory.Jet.FormGap.psi_near_one},
\begin{multline*}
\sum_{a,b}E_{ab}\langle y,F(\cdot)G_{ab}\rangle=\Bigl\langle P_Sy,\sigma(\cdot)\sum_{a,b}E_{ab}G_{ab}\Bigr\rangle\\
-d_S\sum_{a,b}E_{ab}\Bigl(z\mapsto\int_{K_\eps}\Phi_\lambda(\overline{k^{-1}zk})\langle y,\sigma(k)w\rangle\langle\sigma(k)w,
G_{ab}\rangle dk\Bigr).
\end{multline*}
The first term vanishes by Schmid's equation $\sum_{a,b}E_{ab}G_{ab}=0$ (\lref{GlobalTheory.Jet.Schmid.mul_vminus_eq_zero}),
and the second vanishes near $1$ by the corresponding identity for the $K$-average of $\Phi_\lambda$
(\lref{GlobalTheory.Jet.SchmidIntegral.eab_integral_eq_zero}). (d) $F(1)=P_S-\Psi(1)P_S=0$ on $S$.

\emph{Step 4: conclusion.} By \cref{pf:sh-flag:germ}, all jets $E_\omega\langle v,F(\cdot)u\rangle(1)$ vanish for $u,v\in S$.
The function $t\mapsto\langle v,F(\exp(tX))u\rangle$ is real-analytic: its first part by \cref{pf:sh-flag:lowkan} (note
$\mathfrak p_\eps\subseteq\mathfrak g_\eps$) and its second part because the entries of $\Psi$ are analytic and $t\mapsto\exp(tX)$ is
analytic. By \cref{pf:sh-flag:zerojet}, $\langle v,F(\exp X)u\rangle=0$, that is,
$\langle v,\sigma(\exp X)u\rangle=\langle P_Sv,\Psi(\overline{\exp X})P_Su\rangle$. Now $\overline{\exp X}=\exp\bar X$ with
$\bar X\in\mathfrak p_\eps$, so $\Psi(\overline{\exp X})=\psi(\overline{\exp X})$, and $\overline{\exp X}$ is positive definite, hence
in $U_1$. By \eqref{pf:sh-flag:eq:psi} with $z=\exp X$ and the relation
$\Phi_\lambda(\bar z)=\overline{\Phi_\lambda(z)}=\Phi_\lambda(z)$ for positive definite hermitian $z$ (whose leading minors are
positive), we obtain \eqref{pf:sh-flag:eq:compression}.
\end{proof}

\begin{remark}\label{pf:sh-flag:rem:formulagap}
(a) The gap condition $h\ge120nr^{3/2}$ is used only through $g_0>2n$ (Steps 1 and 3); the stronger condition is the one
under which the decay estimates of \ref{pf:rows:S1} are derived from the formula. (b) The vector $w$ is any unit vector
of weight $\Lambda_\lambda$ in $S$; since $S$ is a single $K_\eps$-type, such vectors are unique up to a scalar of modulus one,
which does not affect \eqref{pf:sh-flag:eq:compression}.
\end{remark}

%==============================================================================
\subsection{Covolumes of unipotent groups and the factorisation of amplitudes}\label{sec:pf:sh-flag:cov}

\subsubsection{The covolume of the Bessel unipotent group}

\begin{definition}[The covolume $\shflagcovN$]\label{pf:sh-flag:defcovn}
Let $k,r\ge0$ and let $H_{k,r}$ be the standard hermitian form of signature type of the Bessel data on $\C^{2r+k+1}$, with
the standard isotropic vectors $b_{e(a)}$ and dual vectors $b_{e^*(a)}$ ($0\le a<r$). Let $N_{k,r}\subseteq\Uni(H_{k,r})$
be the unipotent group of the Bessel data: the isometries which act unipotently (upper triangularly with diagonal
$1$) on the standard isotropic flag $b_{e(0)},\dots,b_{e(r-1)}$ and trivially on the $(k+1)$-dimensional middle block
modulo this flag. Let $\mu_N$ be the push-forward under the exponential map of the Lebesgue measure on the real Lie
algebra of $N_{k,r}$ (in its standard real coordinates); it is a Haar measure on $N_{k,r}$
(\lref{GlobalTheory.haarLie_isHaarOn_Nset}).

For a datum $D$ (ranks $k$ and $r=r_S$), a Hida setting $X$, a full system $Y$ over $X$ and a word $u$, the
\emph{Bessel unipotent group} $N_u$ of $u$ is defined over $F$; the full system provides a measurable fundamental set
$\mathcal F_N(u)\subseteq G_\infty$ for the integral points $N_u(\OO)$ acting on the archimedean points (every $g$ has a
unique $s\in N_u(\OO)$ with $sg\in\mathcal F_N(u)$). For $\theta\in\R$, a \emph{choice} $C$ of $X$ at $\theta$ provides for
every $\tau\in\Psi$ a pair of mutually inverse matrices $Q_\tau,Q_\tau^{-1}$ transporting the archimedean points of
$X$ at the place $v(\tau)$ of $F$ below $\tau$ to the Bessel model; let $\mathrm{tr}_C(x)=(Q_\tau^{-1}x_{v(\tau)}Q_\tau)_{\tau\in
\Psi}$ (in the coordinates of $\mathrm{Mat}_{2r+k+1}(\C)$). Put
\[
\mu_{N,\theta}=\max\bigl((2v)^{-rk}\,\theta,0\bigr)\,\mu_N,\qquad
\shflagcovN(C,u)=\mu_{N,\theta}^{\otimes\Psi}\bigl(\mathrm{tr}_C(\mathcal F_N(u))\bigr),\qquad e_N=r^2+\tfrac r2 .
\]
\formal{Formalism.L2.SG.LoopTerm.covN}
\end{definition}

\begin{theorem}[Covolume of the Bessel unipotent group]\label{pf:sh-flag:covn}
For every datum $D$, Hida setting $X$, full system $Y$, $\theta>0$, choice $C$ of $X$ at $\theta$ and word $u$,
\[
\shflagcovN(C,u)=\bigl(\theta\,(2v)^{e_N}\bigr)^{f} .
\]
\formallabel{GlobalTheory.CovNFormulaPf.main}
\end{theorem}

The main ingredient is the following covolume formula, a computation of the volume of $N(\OO)\backslash N(\R)$ in terms
of the discriminant. For a CM extension $E/F$, $n=2r+k+1$, a hermitian $J\in M_n(E)$ and $P\in M_{n\times(r+1)}(E)$, we say
that $P=(e_0,\dots,e_{r-1},z)$ is a \emph{frame} with norm $T\in F$ if its columns are linearly independent with Gram
matrix $(h_J(P_a,P_b))_{a,b}=\mathrm{diag}(0,\dots,0,-T)$; $f=(f_0,\dots,f_{r-1})$ is a \emph{dual tuple} if $h_J(f_a,e_b)=\delta_{ab}$,
$h_J(f_a,z)=0$ and $h_J(f_a,f_b)=0$. A family $A=(A_w)_w$ of matrices indexed by the infinite places $w$ of $E$ is a
\emph{model} of $J$ if $A_w^*H_{k,r}A_w=J^w$ for all $w$, and it is \emph{adapted} to $(P,f,s)$, $s\in\R_{>0}^r$, if
$A_we_a^w=s_ab_{e(a)}$ and $A_wf_a^w=s_a^{-1}b_{e^*(a)}$ for all $w$ and $a$. For an $\OO_E$-lattice $\Lambda$ let
$N(\Lambda)$ be the group of $\gamma\in\GL_n(E)$ in the unipotent group of the frame $P$ with $\gamma\Lambda=\Lambda$; a
measurable set $\mathcal D\subseteq\prod_wN_{k,r}$ is a \emph{fundamental domain} for $N(\Lambda)$ if for every
$x\in\prod_wN_{k,r}$ there is exactly one $\gamma\in N(\Lambda)$ with $(A_w\gamma^wA_w^{-1}x_w)_w\in\mathcal D$.

\begin{proposition}[Covolume of the integral unipotent radical]\label{pf:sh-flag:covnprop}
Let $E/F$ be a CM extension unramified at all finite primes, $J$ hermitian and non-degenerate, $\Lambda$ an
$\OO_E$-lattice which is self-dual for $h_J$, $P$ a frame with norm $T\in F$, $f$ a dual tuple such that the lattice
spanned by the multiples of the columns of $P$ and $f$ by suitable powers of the inverse different of $E$ is contained in
$\Lambda$, $A$ a model of $J$ adapted to $(P,f,s)$. Then fundamental domains for $N(\Lambda)$ exist, and every
fundamental domain $\mathcal D$ satisfies
\[
\mu_N^{\otimes\{w\}}(\mathcal D)=|d_F|^{-\left((k+1)r^2+\frac43(r^3-r)+\frac32r\right)}
\Bigl(\prod_{a=0}^{r-1}s_a^{4(r-a)+2k}\Bigr)^{[F:\Q]} .
\]
\formallabel{GlobalTheory.T.covolume_integral_unipotentRadical_real_points}\formallabel{GlobalTheory.SnMass.VolN.covolume_N}
\end{proposition}

\begin{proof}
Write $f_E$ for the number of infinite places of $E$ ($f_E=[F:\Q]$) and $\mathfrak d_j\subseteq E$ ($j\in\N$) for the fractional ideals
given by the powers of the inverse different which occur in the hypothesis on the lattice spanned by $P$ and $f$.

\emph{Step 1: the group and its action.} The group $G=N(\Lambda)$ consists of units
(\lref{GlobalTheory.T.integralUnipotentRadical_forms_unit_subgroup}), and $\rho(\gamma)x=(A_w\gamma^wA_w^{-1}x_w)_w$ is an action of
$G$ on $\mathcal N=\prod_wN_{k,r}$ which is faithful (\lref{GlobalTheory.T.unipotentModelAction_isAction_and_faithful}); the elements
$A_w\gamma^wA_w^{-1}$ lie in $N_{k,r}$ (\lref{GlobalTheory.T.integralUnipotent_mem_modelUnipotentAtPlace}), and $\rho(\gamma)$
preserves $\mu_N^{\otimes}$ (\lref{GlobalTheory.T.unipotentModelProductAction_measurePreserving}).

\emph{Step 2: block coordinates.} The coordinates of $N_{k,r}$ are grouped into blocks $\beta$ of four kinds: for pairs $a<b$ of
indices of the flag, an \emph{upper} entry and a \emph{corner} entry; for every $a$, a \emph{middle row} with $k+1$ entries and an
\emph{antidiagonal} entry (purely imaginary, by the unitarity). Each block has a degree, and the coordinates are \emph{layered}: if
$\gamma\in G$ has vanishing coordinates in all blocks of degree $<d$, then for every block $\beta$ of degree $\le d$ and every
$x\in\mathcal N$, $c_\beta(\rho(\gamma)x)=c_\beta(x)+c_\beta(\gamma)$
(\lref{GlobalTheory.T.unipotentBlockEntries_mul_of_hasFlagDepthAtLeast}).

\emph{Step 3: block lattices.} For each block there is a lattice $L_\beta$ in the coordinate space of the block (a product over the
places $w$): for an upper pair $(a,b)$ the ideal $\mathfrak d_{r-a}\mathfrak d_{r-b}^{-1}$ scaled by $s_a/s_b$; for a corner pair
the ideal $\mathfrak d_{r-a}\mathfrak d_{r-b}$ scaled by $s_as_b$; for a middle row a lattice of homomorphisms defined by $\Lambda$,
scaled by $s_a$; for an antidiagonal entry a skew lattice with the ideal $\mathfrak d_{2(r-a)}$ scaled by $s_a^2$. If $\gamma\in G$ has
vanishing coordinates in all blocks of degree $<\deg\beta$, then $c_\beta(\gamma)\in L_\beta$
(\lref{GlobalTheory.T.integralUnipotent_entry_mem_inverseDifferentLattice}, \lref{GlobalTheory.T.integralUnipotent_middleRow_mem_homLattice}, \lref{GlobalTheory.T.integralUnipotent_corner_entry_mem_idealLattice}, and \lref{GlobalTheory.T.integralUnipotent_antidiagonal_mem_skewLattice}); conversely, every $l\in L_\beta$ is $c_\beta(\gamma)$ for some
$\gamma\in G$ whose coordinates in the other blocks of degree $\le\deg\beta$ vanish
(\lref{GlobalTheory.T.exists_integral_unipotent_lift_entry}, \lref{GlobalTheory.T.lift_row_middleSpace_integralPoint}, \lref{GlobalTheory.T.exists_integral_unipotent_lift_corner}, and \lref{GlobalTheory.T.exists_integral_unipotent_lift_antidiagonal}). Each
$L_\beta$ has a measurable fundamental domain $D_\beta$ (every vector has a unique translate by $L_\beta$ in $D_\beta$) with Lebesgue volume
\begin{gather*}
\vol D_{(a,b)}^{\mathrm{up}}=2^{-f_E}|d_F|^{1-2(r-a)+2(r-b)}(s_a/s_b)^{2f_E},\\
\vol D_{(a,b)}^{\mathrm{cor}}=2^{-f_E}|d_F|^{1-2(r-a)-2(r-b)}(s_as_b)^{2f_E},\\
\vol D_a^{\mathrm{mid}}=\bigl(2^{-f_E}|d_F|^{1-2(r-a)}s_a^{2f_E}\bigr)^{k+1},\qquad
\vol D_a^{\mathrm{anti}}=|d_F|^{\frac12-2(r-a)}s_a^{2f_E}
\end{gather*}
(\lref{GlobalTheory.T.covolume_inverseDifferentPow_div_lattice}, \lref{GlobalTheory.T.covolume_inverseDifferentPow_mul_lattice}, \lref{GlobalTheory.T.covolume_middleLattice_homLattice_ideal}, and \lref{GlobalTheory.T.skewLattice_inverseDifferentPow_covolume}); these
are the covolume computations, which use that $E/F$ is unramified, so that the different of $E/F$ is trivial.

\emph{Step 4: a fundamental domain.} By the abstract statement on layered actions
(\lref{GlobalTheory.T.existsUnique_layered_strictFundamentalDomain}), Steps 2 and 3 and the faithfulness imply that
$\mathcal D_0=\{x\in\mathcal N:c_\beta(x)\in D_\beta\text{ for all }\beta\}$ meets every $G$-orbit in $\mathcal N$ in exactly one point:
one adjusts the coordinates block by block, in the order of the degrees, using the lifts of Step 3. So $\mathcal D_0$ is a fundamental
domain in the sense above (the two notions agree via $\gamma\mapsto\rho(\gamma)$).

\emph{Step 5: measures.} The set $\mathcal N$ is measurable of full measure for $\mu_N^{\otimes}$ in the ambient space
(\lref{GlobalTheory.T.pi_unipotentModel_measurable_full_measure}); $G$ is countable, and since $\rho$ preserves the measure, any
two measurable sets meeting every orbit exactly once have the same measure
(\lref{GlobalTheory.T.strictFundamentalDomain_measure_eq}). In the block coordinates the measure $\mu_N^{\otimes}$ is
$2^{f_E(r(r-1)+r(k+1))}$ times the product of the Lebesgue measures of the blocks
(\lref{GlobalTheory.T.measure_unipotent_tuples_free_entry_blocks}, from the description of $\mu_N$ in nilpotent coordinates,
\lref{GlobalTheory.T.haar_lie_nilpotent_coordinates}). Hence every fundamental domain $\mathcal D$ has measure
$2^{f_E(r(r-1)+r(k+1))}\prod_\beta\vol D_\beta$. There are $r(r-1)/2$ upper and $r(r-1)/2$ corner pairs, so the powers of $2$ cancel. The
exponent of $|d_F|$ is
\[
\sum_{1\le j<i\le r}\bigl((1-2i+2j)+(1-2i-2j)\bigr)+\sum_{i=1}^r\bigl((k+1)(1-2i)+\tfrac12-2i\bigr)
=-(k+1)r^2-\tfrac43(r^3-r)-\tfrac32r,
\]
and the exponent of $s_a$ is $f_E\bigl(4(r-1-a)+2(k+1)+2\bigr)=f_E\bigl(4(r-a)+2k\bigr)$
(\lref{GlobalTheory.T.unipotentRadical_blockCovolumes_prod}).
\end{proof}

\begin{lemma}[Normalising constants]\label{pf:sh-flag:normconst}
Let $s_a=(2v)^{r-a}$ for $0\le a<r$. Then $s_a>0$ and
\[
|d_F|^{-\left((k+1)r^2+\frac43(r^3-r)+\frac32r\right)}\Bigl(\prod_{a=0}^{r-1}s_a^{4(r-a)+2k}\Bigr)^{f}=(2v)^{(r^2+rk+\frac r2)f} .
\]
\formallabel{GlobalTheory.T.fundamental_domain_normalising_constants}\formallabel{GlobalTheory.CovDom.normalising_constants}
\end{lemma}

\begin{proof}
Since $E/F$ is unramified at all finite primes, $|d_E|=d_F^2$ (\lref{GlobalTheory.cm_discr_sq}), and $[E:\Q]=2f$; hence
$\mathrm{rd}_E^{2f}=|d_E|=d_F^2$ and $|d_F|=\mathrm{rd}_E^{\,f}=(2v)^f$. Moreover $E$ has $f$ infinite places (they correspond
to the elements of the CM type). The product is $(2v)$ to the power $f\sum_{a=0}^{r-1}(r-a)(4(r-a)+2k)$, and
\[
\sum_{i=1}^ri(4i+2k)=\tfrac23r(r+1)(2r+1)+kr(r+1).
\]
The exponent of $2v$ on the left is therefore $f$ times
\[
-(k+1)r^2-\tfrac43(r^3-r)-\tfrac32r+\tfrac23r(r+1)(2r+1)+kr(r+1)=r^2+rk+\tfrac r2,
\]
as $\frac23r(r+1)(2r+1)-\frac43(r^3-r)=2r^2+2r$.
\end{proof}

\begin{lemma}[The transported fundamental set]\label{pf:sh-flag:covdom}
For every datum $D$, Hida setting $X$, full system $Y$, $\theta\in\R$, choice $C$ of $X$ at $\theta$ and word $u$,
\begin{equation}\label{pf:sh-flag:eq:covdom}
\mu_N^{\otimes\Psi}\bigl(\mathrm{tr}_C(\mathcal F_N(u))\bigr)=(2v)^{(r^2+rk+\frac r2)f}.
\end{equation}
\formallabel{GlobalTheory.T.transportedFundDomain_volume_classical}\formallabel{GlobalTheory.CovDom.main}
\formallabel{GlobalTheory.CovDom.isDomN}
\end{lemma}

\begin{proof}
Let $J_c$, $\Lambda_c$ be the Gram matrix and the lattice of $X$ in the coordinates of the development, $P_c(u)$ the frame of
$u$ with norm $T_c(u)$, $f_c$ the dual tuple, and $A_c(w)$ the transport matrices of $C$ at the infinite places $w$ of $E$.
The hypotheses of \cref{pf:sh-flag:covnprop} hold with $s_a=(2v)^{r-a}$: $E/F$ is a CM extension and $\OO_E$ is unramified
over $\OO_F$ (by the third condition of $(\mathrm{CMU})$ for the host); $J_c$ is hermitian and non-degenerate and
$\Lambda_c$ is self-dual (\lref{GlobalTheory.T.concrete_gram_hermitian_nondegenerate_lattice_selfDual}); $P_c(u)$ is a
frame, $f_c$ a dual tuple, and the lattice they span with the powers of the inverse different lies in $\Lambda_c$
(\lref{GlobalTheory.T.wordFrame_dualTuple_gram_integrality}); and $A_c$ is a model of $J_c$ adapted to
$(P_c(u),f_c,s)$ (\lref{GlobalTheory.T.bessel_transport_isometry_frame_adapted}). Let $G_m$ be the identification of
families indexed by $\Psi$ with families indexed by the infinite places of $E$. We claim that
$\mathcal D=G_m(\mathrm{tr}_C(\mathcal F_N(u)))$ is a fundamental domain for $N(\Lambda_c)$. It is measurable
(\lref{GlobalTheory.T.measurableSet_transported_real_points_image}). The map $g\mapsto G_m(\mathrm{tr}_C(g))$ from the
archimedean points of $N_u$ to $\prod_wN_{k,r}$ is bijective and multiplicative
(\lref{GlobalTheory.T.transport_real_points_bijective_multiplicative}), and the integral points correspond: the
elements $\gamma\in N(\Lambda_c)$ are exactly the elements with $A_c(w)\gamma^wA_c(w)^{-1}=G_m(\mathrm{tr}_C(s))_w$ for some
$s\in N_u(\OO)$ (\lref{GlobalTheory.T.integralUnipotent_eq_latticeUnipotent_coords}). Given $x$, write
$x=G_m(\mathrm{tr}_C(g))$ and choose $s\in N_u(\OO)$ with $sg\in\mathcal F_N(u)$; the corresponding $\gamma$ moves $x$ into
$\mathcal D$. If $\gamma'$ also does, with corresponding $s'$, then $s'g\in\mathcal F_N(u)$ by injectivity; since
$N_u(\OO)$ is a group (\lref{Formalism.L2.SG.Pf.IdMeet.noinfEqC_holds}), $s's^{-1}\in N_u(\OO)$ maps $sg$ into the
fundamental set, so $s'=s$; then $A_c(w)\gamma'^wA_c(w)^{-1}=A_c(w)\gamma^wA_c(w)^{-1}$ with $A_c(w)$ invertible (its
determinant is non-zero since $A_c(w)^*H_{k,r}A_c(w)=J_c^w$ and $\det J_c\ne0$), so $\gamma'=\gamma$ by the injectivity of
the embedding $E\to\C$ at $w$. This proves the claim. Since $G_m$ preserves the product measures
(\lref{GlobalTheory.T.measurePi_pointsAtInfinitePlacesCMType_image}), \cref{pf:sh-flag:covnprop} and
\cref{pf:sh-flag:normconst} give \eqref{pf:sh-flag:eq:covdom}.
\end{proof}

\begin{proof}[Proof of \cref{pf:sh-flag:covn}]
Put $t=(2v)^{-rk}\theta>0$ (note $2v=\mathrm{rd}_E\ge1$ and $\theta>0$). Then
$\mu_{N,\theta}=t\mu_N$, and $(t\mu_N)^{\otimes\Psi}=t^f\mu_N^{\otimes\Psi}$ since $|\Psi|=f$. By \cref{pf:sh-flag:covdom},
\[
\shflagcovN(C,u)=\bigl((2v)^{-rk}\theta\bigr)^f(2v)^{(r^2+rk+\frac r2)f}=\bigl(\theta(2v)^{r^2+\frac r2}\bigr)^f .
\qedhere
\]
\end{proof}

\subsubsection{Terms, flags and stabilisers}

We recall the objects attached to the terms of the orbital expansion which enter the remaining results.

\begin{definition}[Flags, Gram matrices, relevant terms]\label{pf:sh-flag:defflags}
For a datum $D$ let $V_D$ be the hermitian space of $D$ over $E$, of dimension $2r+k+1$ ($r=r_S$). A \emph{flag} with norm $T$ is a
family $F=(e_1,\dots,e_r,z)$ in $V_D$ with $e_1,\dots,e_r$ linearly independent, isotropic and mutually orthogonal,
$h(e_a,z)=0$ and $h(z,z)=-T$, $T\ne0$; $X(F)=\langle e_1,\dots,e_r\rangle$. For two flags $F,F'$ the \emph{Gram matrix}
$\widehat Z(F,F')\in M_{r+1}(E)$ has entries $h(F_a,F'_b)$, its leading principal minors are $D_l(F,F')$
($1\le l\le r+1$), and the \emph{weight} is $G(F,F')=D_{r+1}\bar D_{r+1}-TT'D_r\bar D_r$. Every word $u$ has a flag $F_u$ with
norm the level prime $T_u$, and every term $t$ of $\mathcal T(u,u')$ has a row lattice $\Lambda_t$ and a flag $F_t$ with norm
$T_{u'}$; the weight $G(t)$ of the term is (the image of) $G(F_u,F_t)$, and $t$ is \emph{relevant} if the pair $(F_u,F_t)$
is relevant (the generic characters attached to the two flags agree, in the sense of their real parts, on the common elements of
the Lie algebras of the unipotent groups of the two flags). The
development has an abstract term system attached to $D$, and the \emph{concrete data} of a full system $Y$ with finite
spaces provide an identification of it with the term system of $Y$; we use the same letter $t$ for corresponding
terms.
\formal{Formalism.HermitianLatticeCounting.S6.FlagD}\ \formal{Formalism.HermitianLatticeCounting.S6.Zhat}\
\formal{Formalism.HermitianLatticeCounting.S6.Gval}\ \formal{Formalism.HermitianLatticeCounting.S6.Relevant}
\end{definition}

\begin{definition}[Real stabiliser covolume]\label{pf:sh-flag:defvolr}
For a relevant term $t$ with $G(t)\ne0$ which lifts to the term system of $Y$, the development fixes a \emph{pin} of the
lift: a rational point $w$ of $t$, a measure $\mu_R$ on $G_\infty$ normalised at $w$ (the \emph{pinned} measure on the real
stabiliser of $t$), and a measurable fundamental set $K_t$ for the integral stabiliser of $t$ acting on its real
stabiliser. The \emph{real stabiliser covolume} of $t$ is $\shflagvolR(t)=\mu_R(K_t)$, and $\shflagvolR(t)=1$ if $t$ has no
such lift or no pin.
\formal{Formalism.L2.SG.LoopTerm.volRC}
\end{definition}

\begin{definition}[Heights of Gram matrices]\label{pf:sh-flag:defhhz}
For $b>0$, $C\in\R$ and $Z\in M_{r+1}(E)$, with $|x|'_{\tau,l}=|\tau(x)|\,\mathrm{rd}_E^{-2e(r,l)}$, $e(r,l)=\sum_{a=1}^l(r+1-a)$,
put
\begin{equation}\label{pf:sh-flag:eq:hhz}
H_b(C)=b^{Cm^2f},\qquad \mathrm{HHZ}_b(C,Z)=H_b(C)^{Cm}\Bigl(\prod_{\tau\in\Psi}\prod_{l=1}^{r+1}
\max\bigl(1,|\Delta_l(Z)|'_{\tau,l}\,b^{-Cm^2}\bigr)\Bigr)^4,
\end{equation}
where $\Delta_l(Z)$ is the $l$-th leading principal minor.
\formal{Formalism.MainNotions.S.Cnt.Hb}\ \formal{Formalism.MainNotions.S.Cnt.HHZ}
\end{definition}
Note $\mathrm{HHZ}_b(C,Z)\ge1$ for $b\ge1$ and $C\ge0$.

\subsubsection{The covolume of the stabiliser of a relevant term}

\begin{definition}\label{pf:sh-flag:defhyp}
The \emph{counting hypotheses} $\mathrm{Hyp}(D,L_2,u,u';C',C_l,C_g,c_b,C_w,c_T,C_T,C_0,A)$ are those of \cref{pf:stageU:mainrow}
(the standing hypotheses, the standard conditions for $(u,u')$ and further conditions involving $A$).
\formal{Formalism.MainNotions.S.Cnt.Hyp}
\end{definition}

\begin{theorem}[Stabiliser covolumes of relevant terms]\label{pf:sh-flag:stabcov}
For all $C_g,c_b,C_w,c_T,C_T,C_0,A\in\R$ there is $C\in\R$ such that the following holds. Let $D$ be a datum with size data
$L_2$, $X$ a Hida setting, $Y$ a full system with finite spaces and concrete data for $D$ (with Witt data, tie hypotheses
and amplitude data), $u,u'$ words and $C'$, $C_l\ge1$ such that the counting hypotheses
$\mathrm{Hyp}(D,L_2,u,u';C',C_l,C_g,c_b,C_w,c_T,C_T,C_0,A)$ hold. Then for every relevant term $t$ with $G(t)\ne0$,
\[
\shflagvolR(t)\le\mathrm{HHZ}_b\bigl(C,\widehat Z(F_u,F_t)\bigr)^{Cr_S^2},
\]
where $F_u$ is the flag of $u$ and $F_t$ the flag of $t$ given by the Witt data.
\formallabel{GlobalTheory.T.stabiliserCovolume_relevantTerm_le_heightPow}\formallabel{GlobalTheory.LoopVol.volRC_le}
\end{theorem}

The proof uses a Bruhat-type factorisation of the Gram matrix. A \emph{point} $p$ of size $r+1$ over $E$ determines a
matrix $P(p)$ (a monomial pattern with entries in $E$) and a set $\mathrm{roots}(p)$ of positions; the \emph{root
coordinates} of $p$ are the families $c=(c_\tau)_{\tau\in\Psi}$ of complex numbers indexed by $\mathrm{roots}(p)$. The covolume is
bounded by covers in root coordinates as follows. For a term $s$ of the term system of $Y$ we write $R(s)\subseteq G_\infty$ for its
real stabiliser and $\Gamma(s)\subseteq R(s)$ for its integral stabiliser, a countable group.

\begin{lemma}[Fundamental sets and covering sets]\label{pf:sh-flag:packpinned}
Let $s$ be a term of $\mathcal T(u,u')$ in the term system of $Y$ whose real point is $g\gamma g'^{-1}$, with $g$, $g'$ in the
archimedean stabiliser groups of $u$, $u'$. Let $\mu$ be a Haar measure on the stabiliser $\mathrm{St}_\infty(\gamma)$ of $\gamma$
(a measure concentrated on $\mathrm{St}_\infty(\gamma)$ and invariant under left translation by its elements), $\mu_R$ its image under
$y\mapsto gyg^{-1}$, $K$ a fundamental set for $\Gamma(s)$ in $R(s)$ and $B\subseteq G_\infty$ a measurable set such that for every $x\in R(s)$
there is $q\in\Gamma(s)$ with $qx\in B$. Then $\mu_R(K)\le\mu_R(B)$.
\formallabel{GlobalTheory.LoopVol.packing_pinned}
\end{lemma}

\begin{proof}
Write $\Gamma=\Gamma(s)$. The set $R(s)$ is closed, and $R(s)=g\,\mathrm{St}_\infty(\gamma)\,g^{-1}$; hence $\mu_R$ is concentrated on $R(s)$,
and for $q\in\Gamma\subseteq R(s)$ we have $g^{-1}qg\in\mathrm{St}_\infty(\gamma)$, so left translation by $q$ preserves $\mu_R$. By the
covering property and countable subadditivity,
\[
\mu_R(K)=\mu_R(K\cap R(s))\le\sum_{q\in\Gamma}\mu_R(K\cap q^{-1}B)=\sum_{q\in\Gamma}\mu_R(B\cap qK),
\]
using the invariance under $q$. Since $K$ is a fundamental set for $\Gamma$ in $R(s)$, the sets $B\cap qK$, $q\in\Gamma$, partition $B\cap R(s)$,
and $\mu_R(B\smallsetminus R(s))=0$ (\lref{Formalism.L2.SG.Pf.FormAllPf.decomp}); so the sum is $\mu_R(B)$.
\end{proof}

\begin{lemma}[Covolumes and covers]\label{pf:sh-flag:cover}
Assume the counting hypotheses. Let $s$ be a relevant term of the term system of $Y$ with $G(s)\ne0$ and a pin $(w,\mu_R,K)$, let $L$, $U$
be lower and upper unipotent with $\widehat Z=L\,P(w)\,U$ for the Gram matrix $\widehat Z$ of the pair formed by the base flag of $u$ and the
transported frame of $s$, and let $\Omega$ be a measurable set of root coordinates of $w$ which is reached from every family of root
coordinates by integral elements of the common stabiliser, as in the proof of \cref{pf:sh-flag:stabcov} below. Then
\[
\mu_R(K)\le\mathrm{vol}(\Omega)\,\mathrm{rd}_E^{\,2f(r+1)^3}.
\]
\formallabel{GlobalTheory.T.stabiliserCovolume_le_rootCoordinateCover}\formallabel{GlobalTheory.LoopVol.covolume_le_cover}
\end{lemma}

\begin{proof}
By the definition of a pin, the real point of $s$ is $g\gamma_wg'^{-1}$ with $g$, $g'$ in the archimedean stabiliser groups of $u$, $u'$,
and $\mu_R$ is the image under $y\mapsto gyg^{-1}$ of a Haar measure $\mu_w$ on the stabiliser $\mathrm{St}_\infty(\gamma_w)$ of $\gamma_w$. By
\lref{GlobalTheory.T.integerPoints_cover_realPoints_rootCoordinates} there is a measure-preserving map $T$ such that the set
$B$ of points of $G_\infty$ whose rational coordinates at $w$ lie in $T^{-1}(\Omega)$ has the following property: for every $x$ in the real
stabiliser $R(s)$ there is $q\in\Gamma(s)$ with $qx\in B$. By \cref{pf:sh-flag:packpinned}, $\mu_R(K)\le\mu_R(B)$.

Finally $\mu_R(B)\le\mathrm{vol}(T^{-1}\Omega)\,\mathrm{rd}_E^{\,2f(r+1)^3}$ by the comparison of the Haar measure of the stabiliser in root
coordinates with the Lebesgue measure (\lref{GlobalTheory.T.stabiliserHaar_rootCoordinates_le_lebesgue}), and
$\mathrm{vol}(T^{-1}\Omega)=\mathrm{vol}(\Omega)$ since $T$ preserves the measure.
\end{proof}

\begin{proof}[Proof of \cref{pf:sh-flag:stabcov}]
Let $C\ge1$ be the constant of \lref{GlobalTheory.T.covolume_integralUnipotentStabilizer_of_term} for the given
constants. This statement says: \emph{under the counting hypotheses, for every base flag $F_0$ of $u$, every row lattice
$\Lambda$ of $F_0$, every flag $F'$ for $u'$ forming a term with $(\Lambda,F_0)$, every point $p$ and all lower, resp.\ upper,
unipotent $L,U\in M_{r+1}(E)$ with $\widehat Z(F_0,F')=L\,P(p)\,U$ (in the ordering of the development), there is a measurable set
$\Omega$ of root coordinates of $p$ with}
\[
\mathrm{vol}(\Omega)\,\mathrm{rd}_E^{\,2f(r+1)^3}\le\mathrm{HHZ}_b\bigl(C,\widehat Z(F_0,F')\bigr)^{Cr^2},
\]
\emph{such that for every family $c$ of root coordinates there is an integral element $n$ of the common stabiliser
$N_\Lambda(F_0)\cap N_\Lambda(F')$, acting on the frame of $F'$ by a matrix $a$ (with signs $(-1)^{i+j}$), for which the family
$\bigl((UaU^{-1})^\tau\,R_p(c_\tau)\bigr)_\tau$, restricted to the positions $\mathrm{roots}(p)$, lies in $\Omega$} ($R_p(c_\tau)$ being
the matrix with the root coordinates $c_\tau$ at the positions $\mathrm{roots}(p)$).

Let $t$ be relevant with $G(t)\ne0$. Since $\mathrm{HHZ}_b\ge1$ (\lref{Formalism.MainNotions.BridgeLemma₂.rows_one_le_HHZ})
and $Cr^2\ge0$, the bound holds trivially if $\shflagvolR(t)=1$ by definition, that is, if $t$ has no lift $s$ to the term
system of $Y$ or the lift has no pin. Otherwise let $s$ be the lift, $(w,\mu_R,K)$ its pin, so
$\shflagvolR(t)=\mu_R(K)$. The lift $s$ is relevant with $G(s)\ne0$ (the identification of term systems preserves
relevance and weights). Since $s$ lies at the point $w$, the Gram matrix of the pair (base flag of $u$, frame of $s$) lies
in the cell of $w$: there are lower and upper unipotent $L,U$ with $\widehat Z=L\,P(w)\,U$
(\lref{Formalism.L2.SG.TermTieOf.cellQ}); this persists after transporting the frame by the Witt isometry of $u$
(\lref{Formalism.L2.SG.Saturated.zhat_mapL}), which also preserves positivity and the non-vanishing of the weight. The
transported data form a term with the row lattice of $s$ in the sense of the statement above (positive definiteness is
preserved by the isometry, \lref{Formalism.L2.SG.Pf.positive_definite_isomorphism}), which provides $\Omega$. By
\cref{pf:sh-flag:cover},
\[
\shflagvolR(t)=\mu_R(K)\le\mathrm{vol}(\Omega)\,\mathrm{rd}_E^{\,2f(r+1)^3}\le\mathrm{HHZ}_b\bigl(C,\widehat Z(F_u,F'')\bigr)^{Cr^2},
\]
where $F''$ is the transported frame of $s$. Finally $F''$ and the flag $F_t$ of $t$ are related by the double coset relation
of row lattices (\lref{Formalism.L2.SG.Pf.double_coset_isomorphism}), and $\mathrm{HHZ}_b$ is invariant under this relation
(\lref{GlobalTheory.LoopVol.hhz_dcrel}), which gives the assertion.
\end{proof}

\subsubsection{The factorisation of the amplitudes of relevant terms}

\begin{definition}[Local amplitudes]\label{pf:sh-flag:defp1}
Let $X$ be a Hida setting with a choice $C$ at $\theta=1$ and Wu--Bessel data: for every $\tau\in\Psi$ and all words $u,u'$ a
local Bessel system at $\tau$ with an amplitude map $J_\tau(\nu,\gamma)$ (an operator between the local Bessel spaces of $u'$
and $u$, attached to a measure $\nu$ and a real point $\gamma$), and constants $c_u(\tau)\in\R$ and $\mathrm{scl}_{u,u'}(\tau)\in\C$.
For a term $t$ of $\mathcal T(u,u')$ let $J_t$ be its global amplitude (an operator between the archimedean spaces of $u'$
and $u$), $\gamma_\tau(t)$ the transported real point of $t$ at $\tau$, $\nu_\tau(t)$ its $R$-measure (the measure with the
root measure property at $\gamma_\tau(t)$ fixed by the development, or $0$ if there is none) and
\[
P_1(t)=\prod_{\tau\in\Psi}\shflagHS{J_\tau(\nu_\tau(t),\gamma_\tau(t))}^2 .
\]
\formal{Formalism.L2.SG.TermTieOf.JtOf}\ \formal{Formalism.L2.SG.LoopTerm.P1}
\end{definition}

For a relevant term $s$ of the term system of $Y$ with a pin $(w,\mu_R,K)$, let $\gamma=(\gamma_v)_v$ be its real point, $W_u$, $W_{u'}$ the
local Whittaker--Bessel vectors at the infinite places $v$ of $F$, and $J_v(\nu_v)$ the local orbital operator of
$y\mapsto W_u(y)\otimes\overline{W_{u'}(\gamma_v^{-1}y)}$ integrated against a measure $\nu_v$ on the quotient of the local group by the
stabiliser $\mathrm{St}_v(\gamma_v)$ of the real point. The normalised amplitude of $s$ is
$\mathrm{amp}(s)=\kappa\,\chi(s)\,O_s$, where $\kappa=\mathrm{scl}(\mathcal F_N(u),\mathcal F_N(u'))\,\mathrm{vol}_{K,2}(X)$ is the product of the scalar
of the fundamental sets and a volume constant of $X$, $\chi(s)$ is the value of the normalised level character at the finite
part of $s$, and $O_s$ is the orbital operator of $s$ (the integral over a fundamental set of its integral stabiliser).

\begin{lemma}[Amplitude of a pinned lift]\label{pf:sh-flag:amplift}
Let $s$ be a relevant term of the term system of $Y$ with $G(s)\ne0$, $(w,\mu_R,K)$ a pin of $s$, and assume that for every infinite
place $v$ the operators $W_u(y)\otimes\overline{W_{u'}(\gamma_v^{-1}y)}$ are invariant under left translation of $y$ by
$\mathrm{St}_v(\gamma_v)$. Then there are Haar measures $\mu_v$ on $\mathrm{St}_v(\gamma_v)$ and quotient measures $\nu_v$, such that the image
of $\mu_{v(\tau)}$ in the Bessel model has the root measure property at every $\tau$, and
\begin{equation}\label{pf:sh-flag:eq:amplift}
\shflagHS{\mathrm{amp}(s)}^2=\kappa^2\,\mu_R(K)^2\prod_v\shflagHS{J_v(\nu_v)}^2 .
\end{equation}
\formallabel{GlobalTheory.T.normalisedLift_amplitude_eq_covolume_mul_local}\formallabel{GlobalTheory.AmpLift.amp_lift}
\end{lemma}

\begin{proof}
The pinned measure is a product $\mu_R=\bigotimes_v\mu_v$ of Haar measures on the local stabilisers, and the images of the $\mu_v$ in the
Bessel model have the root measure property (\lref{GlobalTheory.T.normalisedMeasure_eq_prod_besselRootMeasures}); there are local
quotient measures $\nu_v$ and a quotient measure $\nu$ of $\bigotimes_v\mu_v$ on $G_\infty$
(\lref{GlobalTheory.T.exists_quotient_measure_unipotent_stabilisers}). Since $s$ is relevant, the archimedean characters of $u$ and
$u'$ agree on the real stabiliser of $s$ (\lref{Formalism.L2.SG.Pf.RelChar.relCharC_holds}); as $\mu_R$ is a Haar measure on the real
stabiliser, $\nu$ a quotient measure and $K$ a fundamental set for the integral stabiliser, unfolding the integral over $K$ gives
$O_s=\mu_R(K)\,O(\nu;\gamma)$, where $O(\nu;\gamma)$ is the real orbital integral of $W_u\otimes\overline{W_{u'}(\gamma^{-1}\cdot)}$ against $\nu$
(\lref{Formalism.L2.SG.Pf.formAllC_holds}). By the invariance hypothesis the Hilbert--Schmidt norm of the real orbital integral
factorises over the places: $\shflagHS{O(\nu;\gamma)}^2=\prod_v\shflagHS{J_v(\nu_v)}^2$
(\lref{GlobalTheory.T.realOrbitalIntegral_hsNorm_eq_prod_local}). Finally $|\chi(s)|=1$
(\lref{Formalism.L2.SG.Pf.FibreCounting.norm_chiN}). Hence
$\shflagHS{\mathrm{amp}(s)}^2=\kappa^2|\chi(s)|^2\mu_R(K)^2\shflagHS{O(\nu;\gamma)}^2$ is the right-hand side of
\eqref{pf:sh-flag:eq:amplift}.
\end{proof}

\begin{theorem}[Local Bessel factorisation]\label{pf:sh-flag:jfact}
Let $D$ be a datum satisfying $\mathrm{GH}(D)$, $X$ a Hida setting, $Y$ a full system with finite spaces and an
identification of the abstract term system of $D$ with that of $Y$, $C$ a choice of $X$ at $1$ with Wu--Bessel data, and
$u,u'$ words. Then for every relevant term $t$ with $G(t)\ne0$,
\[
\shflagHS{J_t}^2\cdot\shflagcovN(C,u)\,\shflagcovN(C,u')=\shflagvolR(t)^2\,P_1(t).
\]
\formallabel{GlobalTheory.T.localBesselJFactor_all_words_of_gap}\formallabel{GlobalTheory.JFactTie.jFactN_all}
\end{theorem}

\begin{proof}
The term $t$ lifts to a term $s$ of the term system of $Y$ (\lref{Formalism.L2.SG.Pf.liftOntoC_of}, using the lattice
isomorphisms of the Witt flags); $s$ is relevant with $G(s)\ne0$. Every such term has a pin
(\lref{GlobalTheory.T.relevantLift_normalisedData_exist}); let $(w,\mu_R,K)$ be the pin fixed by the development, so that
$\shflagvolR(t)=\mu_R(K)$ and $J_t$ is the normalised amplitude $\mathrm{amp}(s)$ of $s$. At every infinite place $v$, the
rank-one operators $W_u(\cdot)\otimes\overline{W_{u'}(\gamma_v^{-1}\,\cdot)}$ built from the local Whittaker--Bessel vectors
are invariant under left translation by the stabiliser of the real point $\gamma_v$ of $s$: the development proves this
from the gap hypothesis (\lref{GlobalTheory.StabInv.stab_inv}, applied to a real relative point of $t$ at $\tau$ given by
\lref{GlobalTheory.T.relevant_term_has_real_relative_point}). By \cref{pf:sh-flag:amplift} there are Haar measures $\mu_v$
on the local stabilisers and quotient measures $\nu_v$, with the root measure property at every $\tau$, such that
$\shflagHS{J_t}^2=\kappa^2\,\mu_R(K)^2\prod_v\shflagHS{J_v(\nu_v)}^2$. For every $\tau$, with $v=v(\tau)$, the local amplitude of the Wu--Bessel data is related to $J_v(\nu_v)$ by
\begin{equation}\label{pf:sh-flag:eq:amploc}
\shflagHS{J_\tau(\nu_\tau(t),\gamma_\tau(t))}^2=\bigl|c_u(\tau)^{-1}\mathrm{scl}_{u,u'}(\tau)\bigr|^2\,\shflagHS{J_v(\nu_v)}^2
\end{equation}
(\lref{GlobalTheory.T.localAmplitude_wittLift_eq_quotientIntegral}), and the constants satisfy
\begin{equation}\label{pf:sh-flag:eq:ampconst}
\prod_{\tau\in\Psi}\bigl|c_u(\tau)^{-1}\mathrm{scl}_{u,u'}(\tau)\bigr|^2=\kappa^2\,\shflagcovN(C,u)\,\shflagcovN(C,u')
\end{equation}
(\lref{GlobalTheory.T.amplitudeConstant_eq_prod_covolumes}, with $\theta=1$). Since $\tau\mapsto v(\tau)$ is a bijection from
$\Psi$ onto the infinite places of $F$, multiplying \eqref{pf:sh-flag:eq:amploc} over $\tau$ and using
\eqref{pf:sh-flag:eq:ampconst} and \eqref{pf:sh-flag:eq:amplift} gives
\begin{align*}
\shflagvolR(t)^2P_1(t)&=\mu_R(K)^2\,\kappa^2\shflagcovN(C,u)\shflagcovN(C,u')\prod_v\shflagHS{J_v(\nu_v)}^2\\
&=\shflagHS{J_t}^2\,\shflagcovN(C,u)\shflagcovN(C,u').\qedhere
\end{align*}
\end{proof}

\begin{remark}\label{pf:sh-flag:rem:jfact}
The development rescales the $R$-measures of a relevant term by
$\theta(t)=\bigl(\shflagvolR(t)^2P_1(t)/\shflagHS{J_t}^2\bigr)^{1/(2f)}$ when $J_t\ne0$ and $P_1(t)>0$. By the theorem and
\cref{pf:sh-flag:covn} with $\theta=1$, $\theta(t)=\bigl(\shflagcovN(C,u)\shflagcovN(C,u')\bigr)^{1/(2f)}=(2v)^{e_N}$ for all such
terms.
\end{remark}

%==============================================================================
\subsection{The type bound}\label{sec:pf:sh-flag:type}

\subsubsection{Statement}

\begin{definition}[Labels and weights of relevant terms]\label{pf:sh-flag:deflabels}
Let $D$ be a datum with size data $L_2$ (scale $b$), $X$ a Hida setting, $Y$ a full system with $(X,Y)$ concretely tied, and
$u$ a word; $r=r_S$. The degenerate relevant terms of $\mathcal T(u,u)$ form the set
$\shflagrel(u)=\{t\in\mathrm{Deg}:t\text{ relevant}\}$. For $t\in\shflagrel(u)$ let:
\begin{itemize}
\item $j(t)=r-\dim_E\bigl(X(F_u)\cap X(F_t)\bigr)$, the \emph{corank} of $t$;
\item $\delta(t)=(\delta_l(t))_{l\ge1}$ with $\delta_l(t)=D_l(F_u,F_t)$ for $l\le j(t)$ and $\delta_l(t)=0$ otherwise, the
\emph{minor labels} of $t$;
\item $\mathrm{wt}(t)=\bigl(w(\Lambda_t)w(\Lambda'_t)\bigr)^{1/2}$, the geometric mean of the row weights of the two row lattices of
$t$;
\item $\omega(t)$, the \emph{index weight}: the infimum, over the normal forms $N$ of the pair $(F_u,F_t)$ adapted to the
different $\mathfrak d$ of $F$ (extended to $E$), of the relative inverse volume $\omega_N(\Lambda_t)$ of the stabiliser of the
row lattice $\Lambda_t$ in the group generated by the common unipotent elements of the two flags, compared with the stabiliser
of the reference lattice of $N$.
\end{itemize}
Let $d_\pi(u)$ be the dimension of the archimedean space of $u$ and $M(u)=\sum_\Lambda w(\Lambda)$ the mass of the rows of $u$
(the sum of the row weights over the row lattices of $u$ in $Y$). For $1\le j\le r$ and $\delta\colon\N\to E$ put, with
$e(r,l)=\sum_{a=1}^l(r+1-a)$, $|x|'_{\tau,l}=|\tau(x)|\,\mathrm{rd}_E^{-2e(r,l)}$ and $H_b(C)=b^{Cm^2f}$:
\begin{gather*}
\Delta(j,\delta)=\Nm\bigl(\delta_j\,\mathfrak d^{-2e(r,j)}\bigr),\qquad
S_H(j,\delta)=\prod_{\tau\in\Psi}\max\Bigl(1,\max_{1\le l\le j,\ \delta_l\ne0}\frac{|\delta_l|'_{\tau,l}}{(2H)^{2r}}\Bigr),\\
a(j,\delta)=\Bigl(\frac Ck\Bigr)^{j(k-1-j)f}\Delta(j,\delta)^{-\frac{k-1-j}2}H_b(C)^{Cr^4}S_{H_1}(j,\delta)^{-C_sk},\\
e(j,\delta)=\Delta(j,\delta)^{\frac k2}H_b(C_h)^{C_hr^4}S_{H_1}(j,\delta)^{C_pk},\qquad H_1=b^{C_Hm},
\end{gather*}
where $\Nm$ is the absolute norm of fractional ideals of $E$; $a$ depends on $(C_H,C_s,C)$ and $e$ on $(C_H,C_h,C_p)$.
\formal{Formalism.L2.SG.KappaOffDiagonal.jT}\ \formal{Formalism.L2.SG.KappaOffDiagonal.labT}\ \formal{GlobalTheory.TypeAsm.omT}\
\formal{Formalism.MainNotions.KappaEndpoints.k2a}\ \formal{Formalism.MainNotions.KappaEndpoints.k2e}
\end{definition}

\begin{definition}[The type condition]\label{pf:sh-flag:typecond}
The \emph{type condition} $\mathrm{Ty}(C_H,C_s,C,C_h,C_p)$ at $(L_2,Y,u)$ is: for all $j',j''\in[1,r]$, all
$\delta',\delta''\colon\N\to E$ and every finite set $S$ of pairs $(t',t'')$ of terms in $\shflagrel(u)$ with the same row and with
$j(t')=j'$, $j(t'')=j''$, $\delta(t')=\delta'$, $\delta(t'')=\delta''$,
\[
\sum_{(t',t'')\in S}\|\mathrm{vec}\,t'\|\,\|\mathrm{vec}\,t''\|\le d_\pi(u)\,M(u)\,a(j',\delta')e(j',\delta')\,a(j'',\delta'')e(j'',\delta'').
\]
\formal{Formalism.L2.SG.KappaOffDiagonal.TypeAt}
\end{definition}

\begin{theorem}[The type bound]\label{pf:sh-flag:typebound}
For all $C_g,c_b,C_w,c_T,C_T,C_0$ with $c_b,c_T>0$ there are $C_H,C_h,C_p\ge0$ such that for every $C_s\ge C_p+4$ there is
$C_\circ$ such that for every $C\ge C_\circ$ there are $C'_0,C_{l,0}$ such that for all $C'\ge C'_0$, $C_l\ge C_{l,0}$, every
datum $D$ with size data $L_2$, every concretely tied $(X,Y)$ and every word $u$: if the ties of $D$ and the $K_2$-hypotheses
$\mathrm{K}_2(C',C_l,C_g,c_b,C_w,c_T,C_T,C_0)$ hold, then $\mathrm{Ty}(C_H,C_s,C,C_h,C_p)$ holds at $(L_2,Y,u)$.
\formallabel{GlobalTheory.rowsPos_typeP}\formallabel{GlobalTheory.T.typeBound_of_largeGap_discreteSeries_bessel}
\formallabel{GlobalTheory.TypeAsm.TopGap.typeBound_of_gap}
\end{theorem}

The bound is the product of a bound for single terms and a weighted count of pairs.

\begin{proposition}[Amplitudes of relevant terms]\label{pf:sh-flag:amp}
Assume the uniqueness of Bessel models in large gap and the Bessel--Whittaker statement (the two components of
\lref{GlobalTheory.T.bessel_multiplicity_one_large_gap_and_whittaker_bessel}). For all $C_g,c_b,C_w,c_T,C_T,C_0$ with
$c_b,c_T>0$ there is $C_H\ge0$ such that for every $C_s\ge0$ there is $C_\circ\ge0$ such that for every $C\ge C_\circ$ there are
thresholds such that, for $C',C_l$ beyond them and all $D,L_2,X,Y,u$ as in \cref{pf:sh-flag:typebound}, every
$t\in\shflagrel(u)$ satisfies
\begin{equation}\label{pf:sh-flag:eq:amp}
\|\mathrm{vec}\,t\|\le\mathrm{wt}(t)\,\sqrt{d_\pi(u)}\;\omega(t)\,a\bigl(j(t),\delta(t)\bigr).
\end{equation}
\formallabel{GlobalTheory.T.amplitudeBound_at_terms_of_localData}
\formallabel{GlobalTheory.TypeAsm.AmpTerms.amplitudeBound_of_localData_of_tailBound}
\end{proposition}

\begin{proposition}[Weighted pairs of relevant terms]\label{pf:sh-flag:pairs}
For all $C_g,c_b,C_w,c_T,C_T,C_0$ and $C_H\ge0$ there are $C_h\ge0$, $C_p=16$ and thresholds such that, for $C',C_l$ beyond them,
every datum $D$ with size data $L_2$, every concretely tied $(X,Y)$ and every word $u$ satisfying the $K_2$-hypotheses, and all
$j',j'',\delta',\delta'',S$ as in \cref{pf:sh-flag:typecond},
\begin{equation}\label{pf:sh-flag:eq:pairs}
\sum_{(t',t'')\in S}\mathrm{wt}(t')\mathrm{wt}(t'')\,\omega(t')\omega(t'')\le M(u)\,e(j',\delta')\,e(j'',\delta'').
\end{equation}
\formallabel{GlobalTheory.T.termPairCountBound_of_siegelMassRatios}
\formallabel{GlobalTheory.TypeAsm.pair_count_bound_of_siegel_mass}
\end{proposition}

\begin{lemma}[Combination]\label{pf:sh-flag:comb}
The conclusions of \cref{pf:sh-flag:amp} (for all constants) and of \cref{pf:sh-flag:pairs} (for all constants) imply the
conclusion of \cref{pf:sh-flag:typebound}.
\formallabel{GlobalTheory.TypeAsm.typeBound_of_amplitude_pairs}
\end{lemma}

\begin{proof}
Given $C_g,c_b,C_w,c_T,C_T,C_0$ with $c_b,c_T>0$, let $C_H\ge0$ be the constant of \cref{pf:sh-flag:amp}, and let $C_h,C_p\ge0$
and thresholds $C'_1,C_{l,1}$ be given by \cref{pf:sh-flag:pairs} for these constants and $C_H$. Let $C_s\ge C_p+4$; then
$C_s\ge0$, and \cref{pf:sh-flag:amp} gives $C_\circ\ge0$; for $C\ge C_\circ$ it gives thresholds $C'_2,C_{l,2}$. Take
$C'_0=\max(C'_1,C'_2)$ and $C_{l,0}=\max(C_{l,1},C_{l,2})$. Let $j',j'',\delta',\delta'',S$ be as in \cref{pf:sh-flag:typecond}. Since
$C\ge0$, the factors $a(j,\delta)$ are non-negative (\lref{GlobalTheory.TypeAsm.termAmplitudeFactor_nonneg}). For
$(t',t'')\in S$, \eqref{pf:sh-flag:eq:amp} for $t'$ and $t''$ gives
\[
\|\mathrm{vec}\,t'\|\,\|\mathrm{vec}\,t''\|\le d_\pi(u)\,a(j',\delta')a(j'',\delta'')\,\mathrm{wt}(t')\mathrm{wt}(t'')\,\omega(t')\omega(t''),
\]
using $\sqrt{d_\pi(u)}^2=d_\pi(u)$. Summing over $S$ and applying \eqref{pf:sh-flag:eq:pairs} (multiplied by the non-negative
factor $d_\pi(u)a(j',\delta')a(j'',\delta'')$) gives the type condition.
\end{proof}

\begin{proof}[Proof of \cref{pf:sh-flag:typebound}]
The uniqueness of Bessel models in large gap and the Bessel--Whittaker statement hold
(\lref{GlobalTheory.T.bessel_multiplicity_one_large_gap_and_whittaker_bessel}); the first gives the gap hypotheses
of the development (\lref{GlobalTheory.GapIn.hyp_of_uniqGap}). Under these, \cref{pf:sh-flag:amp} holds. The maximal
lattice mass formula (\lref{GlobalTheory.Mass.MaxFormula.main}: for every CM field satisfying $(\mathrm{CMU})$ and
every $n\ge1$ there is $c$ such that the mass of every maximal lattice of every positive definite hermitian space of
dimension $n$ is $c$ times a product of local factors) and the Siegel mass ratio (\cref{pf:sh-flag:massratio}) give
\cref{pf:sh-flag:pairs}. \Cref{pf:sh-flag:comb} concludes.
\end{proof}

\subsubsection{The amplitude of a relevant term}

The bound \eqref{pf:sh-flag:eq:amp} is obtained by factorising the amplitude of a term into the covolume of its integral
stabiliser and local amplitudes at the real places, and by bounding the latter by explicit degree bounds. We list the
ingredients; they are statements of the development for all data satisfying the $K_2$-hypotheses (written $\mathrm{Std}$).
\begin{enumerate}[label=(A\arabic*)]
\item\label{pf:sh-flag:A1} \emph{Local data} (\lref{GlobalTheory.TypeAsm.locDatC_of_besselOfW}, from the
Bessel--Whittaker statement): for every datum with $\mathrm{GH}(D)$ and the $d_\sigma$-tie, every Hida setting $X$ has a choice
$C$ at $1$ with Wu--Bessel data which is \emph{good} (local Bessel systems at every $\tau$ with the properties required
below).
\item\label{pf:sh-flag:A2} \emph{Tail bounds} (\lref{GlobalTheory.T.localDataTailBound_closedForm_largeGaps}, from the gap
hypotheses): for all constants with $c_b,c_T>0$ there is $C_H\ge0$ such that for every $C_1\in\N$, beyond thresholds, the
local Bessel data satisfy tail bounds of order $q=C_1kr$ with height $H_1=b^{C_Hm}$.
\item\label{pf:sh-flag:A3} \emph{Product of degree bounds} (\lref{GlobalTheory.T.localBound_productOverPlaces_label}): for
every $C_s\ge0$ there is $C_1\ge1$ such that for $k\ge40r^4$, $k\ge5r$ and every $t\in\shflagrel(u)$ with $1\le j(t)\le r$ there
are $P>0$ with $\Delta(j(t),\delta(t))=P^2$ and local discriminants $D_\tau$ such that the product over $\tau$ of the local
degree bounds $\mathrm{deg}(k,r,j,C_1kr,H_1,D_\tau,\dots)$ is at most
\[
\Bigl(\frac{16\pi e}k\Bigr)^{j(k-1-j)f}P^{-(k-1-j)}\,(8rH_1)^{80r^4f}\,S_{H_1}(j,\delta)^{-C_sk},\qquad j=j(t),\ \delta=\delta(t).
\]
\item\label{pf:sh-flag:A4} \emph{Local amplitude bound} (\lref{GlobalTheory.T.amplitude_hsNorm_le_degreeBound_realPlace},
about $4\,000$ lines): for a good local Bessel system at $\tau$ with tail bounds of order $q\ge16r^3+4r+2$, $1\le j\le r$,
$j+2\le k$, a real point $\gamma$, a normal-form stabiliser measure $\mu_R$ with scale $s>0$ and a quotient measure $\nu_R$,
\[
\shflagHS{J_\tau(\nu_R,\gamma)}\le\sqrt{\dim W_\tau}\;\frac{(2v)^{-rk}\,\theta\,2^{r(k+1)+r(r-1)}}s\;
\mathrm{deg}(k,r,j,q,H_1,D_\tau,\dots),
\]
where $W_\tau$ is the local archimedean space of $u$ at $\tau$.
\item\label{pf:sh-flag:A5} \emph{Factorisation}: \cref{pf:sh-flag:fact} below.
\item\label{pf:sh-flag:A6} \emph{Covolume of the integral stabiliser} (\lref{GlobalTheory.T.integralStabiliser_covolume_le_indexWeight}):
for every normal form $N$ adapted to $(F_u,F_t)$ and normal-form stabiliser measures $\mu_\tau$ (of scale $1$) and every
fundamental set $K$,
\[
\mu^{\otimes\Psi}(\mathrm{tr}(K))\le\omega_N(\Lambda_t)\,v^{k(r-j_N)f}\,\ell_V,\qquad \ell_V=(8v)^{2(r+1)^2f}.
\]
\item\label{pf:sh-flag:A7} \emph{Finiteness} (\lref{GlobalTheory.T.integralUnipotent_covolume_finite_bounded}):
$\shflagcovN(C,u)<\infty$ and $\shflagcovN(C,u)\,\ell_V\ge1$.
\item\label{pf:sh-flag:A8} \emph{Normal forms} (\lref{GlobalTheory.T.exists_adaptedNormalForm_relevantTerm},
\lref{GlobalTheory.T.adaptedNormalForm_normsMinors_eq_label}, \lref{GlobalTheory.T.adaptedNormalForm_gives_real_normal_form},
\lref{GlobalTheory.T.strict_fundamentalSet_integralStabilizer_exists}): every $t\in\shflagrel(u)$ has an adapted normal form;
every adapted normal form $N$ has $j_N=j(t)$ and gives a real normal form of $t$ at every $\tau$ (a scaling $\beta$, unipotent
matrices $n_{1,\tau},n_{2,\tau}$ and further data), and the integral stabiliser of $t$ has a fundamental set; moreover $j(t)\ge1$
(\lref{Formalism.L2.SG.KappaOffDiagonal.JPosC_holds}).
\item\label{pf:sh-flag:A9} \emph{Arithmetic} (\lref{GlobalTheory.amplitude_arithmetic}): for every $C_H\ge0$ there is $C_a$
such that for $C\ge C_a$, $m=2r+2$, $1\le j\le r$, $r+2\le k$, $f\ge1$, $b\ge2$, $1\le v$, $2v\le b$, $P>0$, $\Delta=P^2$ and $S\ge1$,
\begin{multline*}
v^{k(r-j)f}\ell_V^2\bigl((2v)^{-rk}2^{r(k+1)+r(r-1)}\bigr)^f\Bigl(\frac{16\pi e}k\Bigr)^{j(k-1-j)f}P^{-(k-1-j)}(8rH_1)^{80r^4f}S^{-C_sk}\\
\le\Bigl(\frac Ck\Bigr)^{j(k-1-j)f}\Delta^{-\frac{k-1-j}2}H_b(C)^{Cr^4}S^{-C_sk}.
\end{multline*}
\end{enumerate}

\begin{lemma}[Factorisation of the amplitude of a term]\label{pf:sh-flag:fact}
Let $C$ be a good choice of $X$ at $1$ with Wu--Bessel data and assume $\shflagcovN(C,u)<\infty$. Let $t$ be a relevant term of
$\mathcal T(u,u)$, $(\mu_\tau)_\tau$ $\sigma$-finite Haar measures on the stabilisers (in the Bessel model) of the real points
$\gamma_\tau(t)$, $(\nu_\tau)_\tau$ $\sigma$-finite measures, concentrated on the model group, such that $\nu_\tau$ is a quotient measure
of the model Haar measure by the image of $\mu_\tau$ under conjugation by $\gamma_\tau(t)^{-1}$, and $K$ a fundamental set for the
integral stabiliser of $t$. Then
\[
\|\mathrm{vec}\,t\|\,\shflagcovN(C,u)\le\mathrm{wt}(t)\,\mu^{\otimes\Psi}(\mathrm{tr}_C(K))\prod_\tau\shflagHS{J_\tau(\nu_\tau,\gamma_\tau(t))}.
\]
\formallabel{GlobalTheory.T.termAmplitude_le_covolume_localAmplitudes}\formallabel{GlobalTheory.TypeAsm.factC_holds}
\formallabel{GlobalTheory.TypeAsm.factC_of}
\end{lemma}

\begin{proof}
We use three statements of the development. (i) \emph{Pull-back of measures}
(\lref{GlobalTheory.T.localHaarMeasures_pullback_realGroup}): there are a Haar measure $\mu_R$ on the real stabiliser of $t$ and
$\sigma$-finite measures $\nu^\flat_v$ on the local groups such that $\bigotimes_v\nu^\flat_v$ is a quotient measure for $\mu_R$,
$\mu_R(K)=\mu^{\otimes\Psi}(\mathrm{tr}_C(K))$ for every measurable $K$, and the image of $\nu^\flat_{v(\tau)}$ in the Bessel model is $\nu_\tau$.
(ii) \emph{A tensor bound} (\lref{GlobalTheory.T.orbitalIntegralTensorBoundHyp_proof}): the Hilbert--Schmidt norm of the real
orbital integral $O(\bigotimes_v\nu^\flat_v;\gamma)$ of $W_u\otimes\overline{W_u(\gamma^{-1}\cdot)}$ is at most $\prod_v\shflagHS{J_v(\nu^\flat_v)}$.
(iii) \emph{Normalising constants} (\lref{GlobalTheory.T.normalisingConstantsProductHyp_proof}): for a good choice with
$\shflagcovN(C,u)<\infty$, $c_u(\tau)>0$ and $\mathrm{scl}_{u,u}(\tau)=1$ for all $\tau$, and $\kappa\prod_\tau c_u(\tau)\,\shflagcovN(C,u)\le1$,
where $\kappa=\mathrm{scl}(\mathcal F_N(u),\mathcal F_N(u))\,\mathrm{vol}_{K,2}(X)\ge0$.

The vector of $t$ is $\mathrm{wt}(t)$ times the image of the normalised amplitude $\mathrm{amp}(t)=\kappa\,\chi(t)\,O_t$ in the
component of the row of $t$, so $\|\mathrm{vec}\,t\|=\mathrm{wt}(t)\shflagHS{\mathrm{amp}(t)}$. As in the proof of \cref{pf:sh-flag:amplift}
(relevance, \lref{Formalism.L2.SG.Pf.formAllC_holds} and $|\chi(t)|=1$),
$\shflagHS{\mathrm{amp}(t)}=\kappa\,\mu_R(K)\,N$ with $N=\shflagHS{O(\bigotimes_v\nu^\flat_v;\gamma)}$. The local amplitudes satisfy
$\shflagHS{J_\tau(\nu_\tau,\gamma_\tau(t))}=c_u(\tau)^{-1}\shflagHS{J_{v(\tau)}(\nu^\flat_{v(\tau)})}$, by the transport of the local amplitude to
the Bessel model (\lref{Formalism.L2.SG.Pf.LocJEq.locJ_map_eq'}) and $\mathrm{scl}_{u,u}(\tau)=1$. Writing $P_v=\prod_v\shflagHS{J_v(\nu^\flat_v)}$
and $P_c=\prod_\tau c_u(\tau)>0$, and using (ii) and (iii),
\begin{align*}
\|\mathrm{vec}\,t\|\,\shflagcovN(C,u)&=\mathrm{wt}(t)\,\kappa\,\mu_R(K)\,N\,\shflagcovN(C,u)
\le\mathrm{wt}(t)\,\mu_R(K)\,\bigl(P_c^{-1}P_v\bigr)\,\bigl(\kappa\,P_c\,\shflagcovN(C,u)\bigr)\\
&\le\mathrm{wt}(t)\,\mu_R(K)\,P_c^{-1}P_v=\mathrm{wt}(t)\,\mu^{\otimes\Psi}(\mathrm{tr}_C(K))\prod_\tau\shflagHS{J_\tau(\nu_\tau,\gamma_\tau(t))}.\qedhere
\end{align*}
\end{proof}

\begin{proof}[Proof of \cref{pf:sh-flag:amp}]
Let $C_H$ be given by \ref{pf:sh-flag:A2}. Let $C_s\ge0$; let $C_1$ be given by \ref{pf:sh-flag:A3} and $C_a$ by \ref{pf:sh-flag:A9},
and put $C_\circ=\max(C_a,0)$. For $C\ge C_\circ$ let $C'_1,C_{l,1}$ be the thresholds of \ref{pf:sh-flag:A2} for $C_1$
and $C_{l,g}$ the threshold of the gap statement (\lref{GlobalTheory.T.gapC_holds}, which gives $\mathrm{GH}(D)$ under the
$K_2$-hypotheses for $C_l\ge C_{l,g}$), and take the thresholds $\max(C'_1,40)$ for $C'$ and $\max(C_{l,1},C_{l,g},1)$ for $C_l$. Let $D,L_2,X,Y,u$ be as in the statement. The $K_2$-hypotheses give $m=2r+2$, $1\le v$, $2v\le b$,
$f\ge1$, $k\ge40r^4$, $k\ge5r$ and a positivity condition at every word (\lref{GlobalTheory.TypeAsm.AmpTerms.sizes}), and that
the host satisfies $(\mathrm{CMU})$. By \ref{pf:sh-flag:A1} there is a good choice $C$ of $X$ at $1$, and by \ref{pf:sh-flag:A2} and
\ref{pf:sh-flag:A7} the tail bounds hold and $\shflagcovN(C,u)\ell_V\ge1$.

Let $t\in\shflagrel(u)$, $j=j(t)$, $\delta=\delta(t)$, and let $P$, $(D_\tau)$ be given by \ref{pf:sh-flag:A3}. Put
$c=\mathrm{wt}(t)\sqrt{d_\pi(u)}\,a(j,\delta)\ge0$. The index weight is the infimum of $\omega_N(\Lambda_t)$ over the adapted normal forms
$N$, which exist by \ref{pf:sh-flag:A8}, and these values are non-negative
(\lref{GlobalTheory.DegAsm.WeightInf.main}); it therefore suffices to show $\|\mathrm{vec}\,t\|\le c\,\omega_N(\Lambda_t)$ for every
adapted normal form $N$.

Fix $N$. By \ref{pf:sh-flag:A8}, $j_N=j$, and there are a real normal form $(\beta,n_{1,\tau},n_{2,\tau})$ and a fundamental set $K$ for
the integral stabiliser. Let $\mu_\tau$ be the image of the Lebesgue measure of the normal-form coordinates with parameter
$\beta$ under conjugation by $n_{2,\tau}^{-1}$; it is a normal-form stabiliser measure of scale $1$, and there are quotient measures
$\nu_\tau$ (\lref{GlobalTheory.DegVol.StabQuot.exists_sigmaFinite_quot}). Then:
\begin{align*}
\|\mathrm{vec}\,t\|&\le\|\mathrm{vec}\,t\|\,\shflagcovN(C,u)\,\ell_V&&\text{(by \ref{pf:sh-flag:A7})}\\
&\le\mathrm{wt}(t)\,\mu^{\otimes\Psi}(\mathrm{tr}(K))\prod_\tau\shflagHS{J_\tau(\nu_\tau,\gamma_\tau(t))}\;\ell_V&&\text{(by \ref{pf:sh-flag:A5})}\\
&\le\mathrm{wt}(t)\,\omega_N(\Lambda_t)v^{k(r-j)f}\ell_V\cdot\sqrt{d_\pi(u)}\bigl((2v)^{-rk}2^{r(k+1)+r(r-1)}\bigr)^f
\prod_\tau\mathrm{deg}_\tau\cdot\ell_V,
\end{align*}
by \ref{pf:sh-flag:A6} and \ref{pf:sh-flag:A4} (with $\theta=s=1$ and $q=C_1kr\ge16r^3+4r+2$, which holds as $k\ge40r^4$), using
$\prod_\tau\sqrt{\dim W_\tau}=\sqrt{d_\pi(u)}$ and $|\Psi|=f$. By \ref{pf:sh-flag:A3} and \ref{pf:sh-flag:A9} (with $S=S_{H_1}(j,\delta)\ge1$
and $\Delta=\Delta(j,\delta)=P^2$), the last expression is at most $\mathrm{wt}(t)\sqrt{d_\pi(u)}\,\omega_N(\Lambda_t)\,a(j,\delta)=c\,\omega_N(\Lambda_t)$.
\end{proof}

\subsubsection{Weighted pairs of relevant terms}

\begin{proof}[Proof of \cref{pf:sh-flag:pairs}]
Given the constants and $C_H\ge0$, put $C_a=\max\bigl(1,2(C_H+1),C_T\bigr)$. The numerical summation lemma
\lref{GlobalTheory.DegAsm.sum_pairBd_mul_datBd_le} gives $C_h\ge0$ and $C_{l}^\sharp$ such that, under explicit size conditions
on $k,r,m,f,b,v$, the level $g$ and the level prime $T\le b^{C_am}$, and for reals $\Delta,\Delta',D_b,D'_b,S,S'\ge1$ with
\[
\Delta\le H_b(C_a)S^2,\quad\Delta'\le H_b(C_a)S'^2,\quad D_b\le(H_b(C_a)S^2)^r,\quad D'_b\le(H_b(C_a)S'^2)^r,
\]
one has
\begin{equation}\label{pf:sh-flag:eq:numsum}
\sum_{x=1}^{\lfloor\Delta\rfloor}\sum_{x'=1}^{\lfloor\Delta'\rfloor}\beta(x,x')\,\alpha(\Delta,D_b;x)\,\alpha(\Delta',D'_b;x')\le
\bigl(\Delta^{\frac k2}H_b(C_h)^{C_hr^4}S^{16k}\bigr)\bigl(\Delta'^{\frac k2}H_b(C_h)^{C_hr^4}S'^{16k}\bigr),
\end{equation}
where $\beta(x,x')$ is an explicit bound for the genus means of embedded pairs of data with indices $x,x'$ and $\alpha$ an explicit
bound for the data (both depending also on $E,k,g,r,T$). Take $C_p=16$ and the thresholds $1$ and $\max(1,C_l^\sharp)$.

Let $D,L_2,X,Y,u$ satisfy the $K_2$-hypotheses and $(X,Y)$ be concretely tied. The host conditions hold
(\lref{GlobalTheory.T.unitaryHostFieldArithmeticConditions}: $E$ satisfies $(\mathrm{CMU})$, $|\Psi|=f$, $[E:\Q]=2f$, $v\ge1$,
$\Nm\mathfrak d=(2v)^{2f}$, $\dim V_D=2r+k+1$, and the level primes are prime and unramified in $E$), and the $K_2$-hypotheses give
the size conditions of the summation lemma (\lref{GlobalTheory.T.parameter_size_constraints}); also
$H_b(2(C_H+1))\le b^{C_am^2f}$ and $T_u^f$ is the norm of the level prime. If $S=\varnothing$ the bound is trivial, since $M(u)\ge0$ and
$e\ge0$. Otherwise the labels $j',\delta'$ and $j'',\delta''$ are those of relevant terms, and for every relevant term
(\lref{GlobalTheory.T.relevantTerm_norms_ge_one_and_bounded})
\begin{gather*}
1\le\Delta(j,\delta),\qquad 1\le D_b(j,\delta),\\
\Delta(j,\delta)\le H_b\bigl(2(C_H+1)\bigr)S_{H_1}(j,\delta)^2,\qquad
D_b(j,\delta)\le\bigl(H_b(2(C_H+1))S_{H_1}(j,\delta)^2\bigr)^r,
\end{gather*}
where $D_b(j,\delta)$ is the product of the norms of the non-zero normalised minors $\delta_l$, $l\le j$. Hence
\eqref{pf:sh-flag:eq:numsum} applies with $\Delta=\Delta(j',\delta')$, $S=S_{H_1}(j',\delta')$, etc., and its right-hand side is
$e(j',\delta')e(j'',\delta'')$ for $C_p=16$.

It remains to bound the left-hand side of \eqref{pf:sh-flag:eq:pairs} by $M(u)$ times the left-hand side of \eqref{pf:sh-flag:eq:numsum}.
There is a positive definite hermitian subspace $W\subseteq V_D$ of dimension $k$, embedded compatibly with the base flag of $u$
(\lref{GlobalTheory.T.definite_quotient_model_exists}). By \lref{GlobalTheory.T.weightedTermPairs_sum_le_via_rows_genusMeans},
the weighted sum of pairs of relevant terms with the given labels is at most $\sum_\Lambda w(\Lambda)\sum_{x,x'}\alpha(x)\alpha(x')\beta(x,x')$,
provided that for every row lattice $\Lambda$ of $u$ and all $x,x'\ge1$ the mean over the genus of the lattice $\Lambda_W$ induced on $W$
(weighted by the inverse orders of the automorphism groups) of the number of embedded pairs of data with indices $x,x'$ is at most
$\shflagmass(\Lambda_W)\,\beta(x,x')$. The ingredients of that statement are the compatibility of terms over rows, the counting of
pairs of terms over a row lattice through embeddings, the row mean statement (\lref{GlobalTheory.T.rowsMeanC_holds}) and the
invariance of the pair counts under isometries. The genus mean bound is
\lref{GlobalTheory.T.genusMean_embeddedPairs_le_mass_mul_pairBound}, which is deduced from the maximal lattice mass formula and
the Siegel mass ratio (\cref{pf:sh-flag:massratio}), with the host conditions and the construction of base lattices. Combining,
\[
\sum_{(t',t'')\in S}\mathrm{wt}(t')\mathrm{wt}(t'')\omega(t')\omega(t'')\le M(u)\sum_{x,x'}\alpha(x)\alpha(x')\beta(x,x')\le
M(u)\,e(j',\delta')\,e(j'',\delta'').\qedhere
\]
\end{proof}

\begin{remark}\label{pf:sh-flag:rem:type}
(a) The ties of $D$ are used only in \cref{pf:sh-flag:amp} (through the $d_\sigma$-tie in \ref{pf:sh-flag:A1} and in the tail
bounds \ref{pf:sh-flag:A2}); \cref{pf:sh-flag:pairs} does not need them, and its proof does not use the concrete tie of $(X,Y)$
either. (b) The constant $C_p=16$ of the pair bound is absolute; the condition $C_s\ge C_p+4$ in
\cref{pf:sh-flag:typebound} is only used to ensure $C_s\ge0$ in the combination. (c) The Siegel--Weil formula is not used.
\end{remark}

%==============================================================================
\subsection{Intersection systems, small sections and generic degrees}\label{sec:pf:sh-flag:geo}

We use arithmetic varieties in coordinates, Hermitian lines and intersection families as in \cref{def:covolume:hs}: $R=\Z[x_{k,i}]$ is
multigraded by a finite set $\mathcal K$, an arithmetic variety of dimension $d$ is a relevant multihomogeneous prime $P\subset R$ with
$P\cap\Z=0$ and $\dim R/P=d+|\mathcal K|$, a Hermitian line $A$ is a multidegree $a$ with a positive function $h_A$ on the complex
multicone, homogeneous of multidegree $a$, and a section of $A$ on $P$ is a multihomogeneous $s\notin P$ of multidegree $a$. A line is
\emph{smooth} if $\log h_A$ is smooth. For an arithmetic variety $P$ of dimension $e+1$, smooth lines $L_1,\dots,L_e$ and a function
$\phi$ on the multicone, $\int_{P(\C)}\phi\,c_1(L_1)\cdots c_1(L_e)$ denotes the \emph{point integral} of the development: the integral
over the complex points of $P$ of $\phi$ against the product of the first Chern forms $dd^c\log h_{L_j}$, defined through a good linear
projection $\Psi$ of the multicone to $\C^{e+1}$ as the integral over the base of the finite fibre sums of the sheet densities
$\phi\cdot\mathrm{MA}(\log h_{L_1},\dots,\log h_{L_e})$. For a prime $p$ and an arithmetic variety $P$, the components $V$ of $P+(p)$
are counted with multiplicities $\mathrm{mult}_V(P+(p))$, and $\deg(V;a_1,\dots,a_e)$ is the geometric intersection number of the
multidegrees on $V$; $\mathrm{int}_\Q(P_\Q;a_1,\dots,a_e)$ is the intersection number of the development on the generic fibre
$P_\Q$.

The \emph{arithmetic Poincar\'e--Lelong formula} $(\mathrm{PL})$ is the hypothesis \ref{H:pf:deepstr:pl} of \cref{sec:pf:deepstr}, and
\emph{resolution in projective coordinates} $(\mathrm{RPC})$ is \cref{def:pf:deepstr:resk}(a).

\begin{definition}\label{pf:sh-flag:defre}
$(\mathrm{RE})$ is the existence of analytic log-resolutions in coordinates: for every arithmetic variety $P$ of dimension $e+1$
and every multihomogeneous ideal $J\not\subseteq P$ there is a multihomogeneous ideal $J'\subseteq J$, $J'\not\subseteq P$, with an analytic
log-resolution of $(P,J')$.
\formal{ReciprocityTheory.ChartRestriction.G.Geo.ResExistsStatement}
\end{definition}
The development derives $(\mathrm{RE})$ from $(\mathrm{RPC})$ (\lref{GlobalTheory.T.resolution_exists_of_resolution_statement}).

\subsubsection{Full intersection systems}

\begin{definition}[Full intersection systems]\label{pf:sh-flag:deffull}
An \emph{intersection system} assigns to every finite index datum $(\mathcal K,(\iota_k))$, arithmetic variety $P$, $n\in\N$ and Hermitian
lines $L_1,\dots,L_n$ a real number $\shflagII(P;L_1,\dots,L_n)$. It is \emph{full} if
\begin{enumerate}[label=(\roman*)]
\item for every index datum, $\shflagII$ is an intersection family (\cref{def:covolume:hs}(c));
\item it is additive in each argument: $\shflagII(P;\dots,A\otimes B,\dots)=\shflagII(P;\dots,A,\dots)+\shflagII(P;\dots,B,\dots)$ for smooth
lines;
\item (projection formula in coordinates) for a block map $(\kappa,v)$ of index data, a point $y$ over a field of characteristic $0$
such that the arithmetic variety $P_y$ of $y$ and its image $v(P_y)$ both have dimension $d$, smooth lines $L_j$ and their pull-backs
$L'_j$: $\shflagII(P_y;L')=[\Q(y):\Q(y\circ v)]\cdot\shflagII(v(P_y);L)$, the degree of the extension of the block function fields;
\item (expected vanishing) if $P'$ has dimension $d+r$ and its image under a block map has dimension $d_0<d$, then
$\shflagII(P';L'_1,\dots,L'_d,N_1,\dots,N_r)=0$ for pull-backs $L'_j$ of smooth lines and smooth $N_i$.
\end{enumerate}
\formal{ReciprocityTheory.ChartRestriction.G.Geo.IsFullIntersectionSystem}
\end{definition}

\begin{theorem}[Existence of a full intersection system]\label{pf:sh-flag:ais}
$(\mathrm{PL})$ and $(\mathrm{RPC})$ imply that a full intersection system exists.
\formallabel{GlobalTheory.CompanionTheory.Ar.Pub.arithIntersectionSystem_of_poincareLelong_resK}
\end{theorem}

The intersection numbers are those of Bost, Gillet and Soul\'e \cite{BostGilletSoule1994,GilletSoule1990}, defined by the recursion of
\cref{def:covolume:hs}(c); the work consists in showing that the recursion is independent of the chosen sections.

\begin{proposition}[Intersection families exist]\label{pf:sh-flag:iexists}
$(\mathrm{PL})$ and $(\mathrm{RE})$ imply that there is an intersection system satisfying (i).
\formallabel{GlobalTheory.T.intersection_exists_of_poincare_lelong}\formallabel{GlobalTheory.CompanionTheory.Ar.Pub.thmIExists}
\end{proposition}

\begin{proof}
The proof establishes the following statements in turn; each is a theorem of the development with the indicated inputs.
\begin{enumerate}[label=(\arabic*)]
\item \emph{Integrability of point integrals}: for an arithmetic variety, a good linear projection, smooth lines and a section $s$
of a smooth line $A$, the fibre sums of the sheet densities of $\log\|s\|_A$ are integrable (from $(\mathrm{RE})$; this is
\cref{pf:sh-flag:ptsint} below), and their absolute values are integrable
(\lref{GroundTheory.Heights₂.Heights.Inter.PtsAbsIntegrable.proof}, with the measurability of the fibre sums).
\item \emph{Bilinearity} of point integrals in the potentials (\lref{GroundTheory.Heights₂.Heights.PtsBilinear.proof}).
\item \emph{Stokes}: for an arithmetic variety of dimension $n+2$, smooth lines $L_1,\dots,L_n$ and smooth lines $B,B'$ of multidegree
$0$, $\int\log h_B\,dd^c\log h_{B'}\wedge c_1(L_1)\cdots c_1(L_n)=\int\log h_{B'}\,dd^c\log h_B\wedge c_1(L_1)\cdots c_1(L_n)$
(\lref{GroundTheory.Heights₂.Heights.StokesStatement.of_found}, from (2), a Segre partition, a Segre chart and Stokes'
theorem on varieties).
\item \emph{An inequality for semipositive lines} (\lref{GroundTheory.Heights₂.Heights.Inter.DeLeSemiPos.proof}), from $(\mathrm{PL})$,
(3), additivity in the first slot, a regularisation limit, monotonicity for semipositive lines, a curvature bound for the
regularisation and the fact that horizontal components do not contain the section.
\item \emph{Symmetry of the section space} (\lref{GroundTheory.Heights₂.Heights.Inter.ProjectiveLineTwo.proof}): for sections
$s,s'$ of smooth lines $A,A'$ on an arithmetic variety of dimension $n+2$ whose divisors have no common horizontal component,
\begin{multline*}
\sum_W\mathrm{mult}_W\!\int_{W(\C)}\log\|s'\|\,c_1(L)+\int_{P(\C)}\log\|s\|\,c_1(A')c_1(L)\\
=\sum_{W'}\mathrm{mult}_{W'}\!\int_{W'(\C)}\log\|s\|\,c_1(L)+\int_{P(\C)}\log\|s'\|\,c_1(A)c_1(L),
\end{multline*}
the sums over the horizontal components of the divisors of $s$, resp.\ $s'$, and $c_1(L)=c_1(L_1)\cdots c_1(L_n)$; from additivity in
the slots, a twist by semipositive lines and (4).
\item \emph{Swap} and \emph{position additivity} (\lref{GroundTheory.Heights₂.Heights.Inter.Swap.proof},
\lref{GroundTheory.Heights₂.Heights.Inter.RPositionAdd.of_remaining}): the recursion may be applied to two lines in either order,
from (5), the cycle relations and a cut of geometric numbers.
\item \emph{Symmetry} of the recursive numbers (\lref{GroundTheory.Heights₂.Heights.Inter.ThmISymm.proof}), from (6), a
comparison of fibres, and the existence, uniqueness and symmetry of intersection numbers over the finite fields $\Z/p$.
\item \emph{Existence} (\lref{GroundTheory.Heights₂.Heights.Inter.ThmIExists.proof}): from the independence of the regularity
order (\lref{GroundTheory.Heights₂.Heights.AutoClose.Inter_R1Indep}), multiplicativity, well-definedness (from (6)) and (7),
the recursion defines an intersection family for every index datum.\qedhere
\end{enumerate}
\end{proof}

\begin{proof}[Proof of \cref{pf:sh-flag:ais}]
By \lref{GlobalTheory.T.resolution_exists_of_resolution_statement}, $(\mathrm{RPC})$ implies $(\mathrm{RE})$, so
\cref{pf:sh-flag:iexists} gives an intersection system $\shflagII$ satisfying (i). Every intersection family is additive in each
argument (\lref{GroundTheory.Heights₂.Heights.Targets.ArithIntersectionSystem.multilin_of_isFamily}), which is (ii). Every intersection
family satisfies the projection formula in coordinates (iii): this is the second part of
\lref{GlobalTheory.T.degree_and_coordinates_of_resolution}, deduced from $(\mathrm{RE})$. Every intersection family satisfies the
expected vanishing (iv) (\lref{GlobalTheory.CompanionTheory.Ar.Pub.contractMixed_closed}: this follows from the vanishing of the
top density of pulled-back forms, the verticality of projective points, the vanishing of the point integrals of forms pulled back
from a variety of smaller dimension, and the arithmetic closure of images). Hence $\shflagII$ is full
(\lref{GroundTheory.Heights₂.Heights.Targets.ArithIntersectionSystem.of_found}).
\end{proof}

\subsubsection{Point integrals and generic degrees}

\begin{proposition}[Integrability of point integrals]\label{pf:sh-flag:ptsint}
Assume $(\mathrm{RPC})$. Let $P$ be an arithmetic variety of dimension $e+1$, $\Psi$ a good linear projection of its multicone to
$\C^{e+1}$, $L_1,\dots,L_e$ and $A$ smooth Hermitian lines and $s$ a section of $A$ on $P$. Then the function
$y\mapsto\sum_{w\in\Psi^{-1}(y)\cap P(\C)}\rho_w(y)$, the finite sum over the fibre of the sheet densities of $\log\|s\|_A$ against
$c_1(L_1)\cdots c_1(L_e)$, is integrable on $\C^{e+1}$.
\formallabel{GlobalTheory.T.points_integrable_of_resolution_statement}
\formal{GlobalTheory.CompanionTheory.Ar.Giv.ptsIntegrable_of_resK}
\end{proposition}

\begin{proof}
$(\mathrm{RPC})$ implies $(\mathrm{RE})$ (\lref{GlobalTheory.T.resolution_exists_of_resolution_statement}). Given the data, $(\mathrm{RE})$
provides an analytic log-resolution $R$ of $P$ (\lref{ReciprocityTheory.ChartRestriction.G.Geo.exists_anLogRes}). On $R$ the
potentials of the lines $L_j$ pull back to potentials, and the function $\log\|s\|$ (lifted to $R$) is integrable against the product of
the pulled-back Chern forms (\lref{ReciprocityTheory.ChartRestriction.G.Geo.brIntegrable}). Its push-forward to the base of $\Psi$ is
integrable (\lref{GroundTheory.Heights₂.Heights.ForSecondHalf.pushForward_integrable}, using that the bad set of the projection is
null and the fibres are finite), and it agrees almost everywhere with the fibre sum of the statement. This is
\lref{GroundTheory.Heights₂.Heights.ForSecondHalf.PtsIntegrable.of_resExists}.
\end{proof}

\begin{proposition}[Chern-form integrals and generic intersection numbers]\label{pf:sh-flag:deggen}
Assume the degree statement: for every arithmetic variety $P$ of dimension $e+1$, smooth lines $L_1,\dots,L_e$ of multidegrees
$a_1,\dots,a_e$ and every prime $p$,
\[
\int_{P(\C)}c_1(L_1)\cdots c_1(L_e)=\sum_{V}\mathrm{mult}_V\bigl(P+(p)\bigr)\deg(V;a_1,\dots,a_e),
\]
the sum over the components $V$ of $P+(p)$. Then for every such $P$ and $L_1,\dots,L_e$,
\[
\int_{P(\C)}c_1(L_1)\cdots c_1(L_e)=\mathrm{int}_\Q(P_\Q;a_1,\dots,a_e).
\]
\formallabel{GlobalTheory.CompanionTheory.Ar.Giv.degreeGeneric_of_degree}
\end{proposition}

\begin{proof}
Apply the degree statement with $p=2$. By the conservation of intersection numbers in the flat family $P$ over $\Z$
(\lref{GroundTheory.Heights₂.Heights.Targets.FibreDegreeSum.proof}: for every arithmetic variety $P$ of dimension $e+1$, all
multidegrees and every prime $p$, $\mathrm{int}_\Q(P_\Q;a_1,\dots,a_e)=\sum_V\mathrm{mult}_V(P+(p))\deg(V;a_1,\dots,a_e)$), the right-hand side
for $p=2$ equals $\mathrm{int}_\Q(P_\Q;a_1,\dots,a_e)$. (The sums are finite, as $P+(p)$ has finitely many minimal primes.)
\end{proof}

\subsubsection{Packing of small sections}

\begin{definition}[Chart lattices and small sections]\label{pf:sh-flag:defsmall}
Let $y$ be a point of $\prod_k\mathbb P^{\iota_k}$ over a field $\Omega$, with arithmetic variety $P_y$ (the kernel of the evaluation at
the multicone of $y$), and assume that $y$ is \emph{birational} (every element of $\Omega$ is a quotient of multihomogeneous polynomials
of the same multidegree evaluated at $y$). For a multidegree $a$ and $\tau\in\N$ let $\Lambda'_\tau(y,a)\subseteq\Omega$ be the group of
elements $s$ such that, for every chart $c$ (a choice of a coordinate $c_k\in\iota_k$ with $y_{k,c_k}\ne0$ for every $k$),
$s\prod_ky_{k,c_k}^{-\tau a_k}$ is integral over the chart ring, the subring of $\Omega$ generated by the ratios $y_{k,i}/y_{k,c_k}$.
For a Hermitian line $A$ of
multidegree $a$, an element $s\in\Omega$ is \emph{small of degree $\tau$} if $|g(z)|\le|h(z)|\,h_A(z)^\tau$ at every complex point $z$ of
the multicone of $P_y$, for all multihomogeneous $g,h$ with $\deg g=\deg h+\tau a$ and $g(y)=s\,h(y)$.
\formal{ReciprocityTheory.ChartRestriction.G.Geo.Birational}\ \formal{ReciprocityTheory.ChartRestriction.G.Geo.chartLatNm}\
\formal{ReciprocityTheory.ChartRestriction.G.Geo.IsSmallH}
\end{definition}

\begin{theorem}[Packing of small sections]\label{pf:sh-flag:packing}
Let $y$ be birational with $P_y$ an arithmetic variety of dimension $D+1$, $A$ a smooth Hermitian line, and $\Lambda_\tau\subseteq
\Lambda'_\tau(y,a)$ subgroups ($\tau\in\N$). Then for every $\tau$ the set of small elements of degree $\tau$ in $\Lambda_\tau$ is finite, and
there is $C\in\R$ such that for all $\tau$
\[
\log\#\{s\in\Lambda_\tau:\ s\text{ small of degree }\tau\}\le C(\tau+1)^{D+1}.
\]
\formallabel{GlobalTheory.T.packing'}\formallabel{GlobalTheory.CompanionTheory.Ar.Small.packing}
\formallabel{MultiProjectiveCoordinates.packing_statement}
\end{theorem}

We use three ingredients.

\begin{lemma}[Packing in normed spaces]\label{pf:sh-flag:packlemma}
Let $V$ be a real normed space, $L\subseteq V$ a finitely generated subgroup, $b\ge0$, and $S\subseteq L$ a finite set of vectors of norm
at most $1$ with $\|s-t\|\ge e^{-b}$ for $s\ne t$ in $S$. Then $\log\#S\le\rank_\Z(L)\,(b+\log3)$.
\formal{GroundTheory.Heights₂.Heights.HilbertSamuel.Packing.proof}
\end{lemma}

This is the usual volume argument in the real span of $L$, whose dimension is at most $\rank_\Z(L)$
(\lref{GroundTheory.Heights₂.Heights.HilbertSamuel.Packing.proof}): the balls of radius $e^{-b}/2$ around the points of $S$ are
disjoint and contained in the ball of radius $1+e^{-b}/2$, so $\#S\le(1+2e^b)^{\dim}\le(3e^b)^{\dim}$.

\begin{lemma}[Rank of the chart lattices]\label{pf:sh-flag:rank}
Under the hypotheses of \cref{pf:sh-flag:packing}, every $\Lambda'_\tau(y,a)$ is a finitely generated free $\Z$-module, and there is
$C_1\ge0$ with $\rank_\Z\Lambda'_\tau(y,a)\le C_1(\tau+1)^D$ for all $\tau$.
\formal{GroundTheory.Heights₂.Heights.HilbertSamuel.NmSectionsRankBound.proof}
\end{lemma}

This is \lref{GroundTheory.Heights₂.Heights.HilbertSamuel.NmSectionsRankBound.proof}, deduced from the finiteness of the chart lattices
(via a conductor of the integral closure) and a bound for the ranks of the groups of sections of multidegree $\tau a$.

\begin{proposition}[Separation of small sections]\label{pf:sh-flag:sep}
Under the hypotheses of \cref{pf:sh-flag:packing} there is $C\ge0$ such that for every $\tau$ there are a finite set $Z$ of complex points
of the multicone of $P_y$ and an additive map $\mathrm{ev}\colon\Lambda'_\tau(y,a)\to\C^Z$ with the following properties: for every
$s\in\Lambda'_\tau(y,a)$ and $z\in Z$ there are multihomogeneous $g,h$ with $\deg g=\deg h+\tau a$, $g(y)=s\,h(y)$, $h(z)\ne0$ and
$\mathrm{ev}(s)(z)\,h(z)=g(z)$; and for all small $s\ne t$ of degree $\tau$ there is $z\in Z$ with
$|\mathrm{ev}(s)(z)-\mathrm{ev}(t)(z)|\ge e^{-C(\tau+1)}h_A(z)^\tau$.
\formallabel{GroundTheory.Heights₂.Heights.HilbertSamuel.SmallSectionsSeparated.proof}
\end{proposition}

\begin{proof}
By the grid separation theorem \lref{MultiProjectiveCoordinates.grid_separation} (about $5\,400$ lines) there is $C\ge0$ such that for
every $\tau$ there is a finite set $Z$ of points of the multicone such that every $s\in\Lambda'_\tau(y,a)$ has, at every $z\in Z$, a
representation $s=g(y)/h(y)$ with $h(z)\ne0$, and every non-zero $u\in\Lambda'_\tau(y,a)$ has, at some $z\in Z$, a representation $u=g(y)/h(y)$
with $h(z)\ne0$ and $|g(z)|\ge e^{-C(\tau+1)}h_A(z)^\tau|h(z)|$. The value $g(z)/h(z)$ does not depend on the representation (two
representations $g/h$, $g'/h'$ of $s$ give $gh'-g'h\in P_y$, which vanishes at the points of the multicone), so
$\mathrm{ev}(s)(z)=g(z)/h(z)$ defines an additive map with the first property. For small $s\ne t$, apply the second property to $u=s-t\ne0$:
at some $z$, $|\mathrm{ev}(s)(z)-\mathrm{ev}(t)(z)|=|\mathrm{ev}(u)(z)|\ge e^{-C(\tau+1)}h_A(z)^\tau$.
\end{proof}

The grid separation theorem is the arithmetic heart of the packing bound: an integral section which is non-zero cannot be small at all
points of a suitable finite grid. Its proof combines representatives on the normal locus, a bound for fibre norms (obtained from
projections avoiding a form, the arithmetic of Segre embeddings, representatives in the fibres, a maximum principle for fibre norms
proved by integrating over the fibres of a pull-back, and a comparison of metrics), and a lower bound on grids of roots; we do not reproduce
it.

\begin{proof}[Proof of \cref{pf:sh-flag:packing}]
Let $C_1$ be given by \cref{pf:sh-flag:rank} and $C$ by \cref{pf:sh-flag:sep}. Fix $\tau$ and let $Z,\mathrm{ev}$ be as in
\cref{pf:sh-flag:sep}. Since $A$ is smooth, $h_A(z)>0$ on the multicone. Consider the additive map
$\phi\colon\Lambda'_\tau(y,a)\to\C^Z$, $\phi(s)(z)=\mathrm{ev}(s)(z)/h_A(z)^\tau$, with the supremum norm on $\C^Z$.

If $s$ is small of degree $\tau$, then $\|\phi(s)\|\le1$: for $z\in Z$ choose $g,h$ as in \cref{pf:sh-flag:sep}; smallness gives
$|g(z)|\le|h(z)|h_A(z)^\tau$, and $g(z)=\mathrm{ev}(s)(z)h(z)$ with $h(z)\ne0$, so $|\mathrm{ev}(s)(z)|\le h_A(z)^\tau$. If $s\ne t$ are small, then
$\|\phi(s)-\phi(t)\|\ge e^{-C(\tau+1)}$ by \cref{pf:sh-flag:sep}; in particular $\phi$ is injective on small elements.

Let $T\subseteq\Lambda_\tau$ be a finite set of small elements. Its image $\phi(T)$ lies in the finitely generated group
$\phi(\Lambda'_\tau(y,a))$, whose rank is at most $\rank_\Z\Lambda'_\tau(y,a)\le C_1(\tau+1)^D$, consists of vectors of norm $\le1$ and is
$e^{-C(\tau+1)}$-separated. By \cref{pf:sh-flag:packlemma},
\begin{equation}\label{pf:sh-flag:eq:pack}
\log\#T\le C_1(\tau+1)^D\bigl(C(\tau+1)+\log3\bigr).
\end{equation}
If the set of small elements of degree $\tau$ in $\Lambda_\tau$ were infinite, it would contain a finite subset $T$ with
$\#T>\exp\bigl(C_1(\tau+1)^D(C(\tau+1)+\log3)\bigr)$, contradicting \eqref{pf:sh-flag:eq:pack}; so it is finite, and
\eqref{pf:sh-flag:eq:pack} applies to it. Since $\tau+1\ge1$ and $\log3\ge0$,
\[
C_1(\tau+1)^D\bigl(C(\tau+1)+\log3\bigr)\le C_1(C+\log3)(\tau+1)^{D+1},
\]
which is the assertion with the constant $C_1(C+\log3)$.
\end{proof}

\begin{remark}\label{pf:sh-flag:rem:packing}
The theorem is used in \cref{lem:pf:deepstr:vi} to compare arithmetic volumes computed with the groups $\Lambda'_\tau$ (elements integral
over the chart rings) and with their subgroups of elements lying in the chart rings; the comparison only needs the finiteness and the
bound for subgroups, which is why the statement allows arbitrary subgroups $\Lambda_\tau\subseteq\Lambda'_\tau(y,a)$.
\end{remark}
